
\documentclass{article}
\usepackage{colortbl}
\usepackage{makecell}
\usepackage{multirow}
\usepackage{supertabular}

\begin{document}

\newcounter{utterance}

\centering \large Interaction Transcript for game `clean\_up', experiment `2\_hard\_7obj\_en', episode 1 with llama{-}3.3{-}70b{-}instruct{-}t1.0.
\vspace{24pt}

{ \footnotesize  \setcounter{utterance}{1}
\setlength{\tabcolsep}{0pt}
\begin{supertabular}{c@{$\;$}|p{.15\linewidth}@{}p{.15\linewidth}p{.15\linewidth}p{.15\linewidth}p{.15\linewidth}p{.15\linewidth}}
    \# & $\;$A & \multicolumn{4}{c}{Game Master} & $\;\:$B\\
    \hline

    \theutterance \stepcounter{utterance}  
    & & \multicolumn{4}{p{0.6\linewidth}}{
        \cellcolor[rgb]{0.9,0.9,0.9}{
            \makecell[{{p{\linewidth}}}]{
                \texttt{\tiny{[P1$\langle$GM]}}
                \texttt{Sono il tuo game master, e stai giocando un gioco collaborativo con la seguente griglia come tavolo da gioco:} \\
\\ 
\texttt{```} \\
\texttt{ 1234567} \\
\texttt{╔═══╤═╤═╗} \\
\texttt{║◌S◌│◌│A║ 1} \\
\texttt{╟─┬─┼─┴─╢ 2} \\
\texttt{║◌│◌│◌◌◌║ 3} \\
\texttt{║P├─┤D┌─╢ 4} \\
\texttt{╟─┘◌│◌│M║ 5} \\
\texttt{║◌◌◌└┬┴─╢ 6} \\
\texttt{║◌OT◌│◌◌║ 7} \\
\texttt{╚════╧══╝} \\
\\ 
\texttt{```} \\
\\ 
\texttt{* Il bordo superiore mostra le coordinate x che aumentano verso destra, e il bordo destro le coordinate y che aumentano verso il basso.} \\
\texttt{* The following objects are randomly placed on your grid: 'P', 'O', 'T', 'S', 'D', 'A', 'M'.} \\
\\ 
\\ 
\texttt{L'altro giocatore vede una variante del tavolo da gioco, dove gli oggetti sono posizionati in diverse posizioni casuali. Non puoi vedere il tavolo dell'altro giocatore, e loro non possono vedere il tuo.} \\
\\ 
\\ 
\texttt{**Obiettivo:**} \\
\\ 
\texttt{Entrambi i giocatori devono spostare gli oggetti sui rispettivi sfondi in modo che oggetti identici finiscano alle stesse coordinate. Devi comunicare con l'altro giocatore per concordare una configurazione obiettivo comune.} \\
\\ 
\\ 
\texttt{**Regole:**} \\
\\ 
\texttt{* In ogni turno, puoi inviare esattamente uno dei seguenti due comandi:} \\
\texttt{ 1. `COMUNICA: <MESSAGE>`: per inviare un messaggio (tutto fino alla prossima interruzione di riga) all'altro giocatore. Lo inoltrerò al tuo partner.} \\
\texttt{2. `SPOSTA: <OBJECT>, (<X>, <Y>)`: per spostare un oggetto in una nuova posizione, dove <X> è la colonna e <Y> è la riga. Ti farò sapere se la tua mossa è andata a buon fine o no.} \\
\texttt{* Se non rispetti il formato, o invii diversi comandi contemporaneamente, devo penalizzarti.} \\
\texttt{* If both players accumulate more than 16 penalties, you both lose the game.} \\
\texttt{* È vitale che comunichi con l'altro giocatore riguardo al tuo stato obiettivo! L'*unico* modo in cui puoi trasmettere la tua strategia all'altro giocatore è usando il comando `COMUNICA: <MESSAGE>`!} \\
\\ 
\\ 
\texttt{**Spostamento Oggetti**} \\
\\ 
\texttt{* Puoi spostare oggetti solo in celle entro i limiti della griglia. La cella di destinazione deve essere vuota, cioè deve contenere solo il simbolo '◌'.} \\
\texttt{* Se provi a spostare un oggetto in un posto che non è vuoto, o provi a spostarlo fuori dalla griglia, devo penalizzarti. Hai un altro tentativo.} \\
\texttt{* Prima di fare una mossa, controlla due volte che il posto di destinazione sia vuoto, e non contenga alcuna lettera, cornice, o linea!} \\
\\ 
\\ 
\texttt{**Fine del Gioco**} \\
\\ 
\texttt{Se pensi di aver raggiunto l'obiettivo di allineare tutti gli oggetti, puoi chiedere all'altro giocatore di finire il gioco inviando COMUNICA: finito?. Se l'altro giocatore ti chiede di finire il gioco, e rispondi con COMUNICA: finito!, the game will end.} \\
\\ 
\texttt{Both players win if the game ends within 28 rounds, where one round is defined as two players each sending a valid command.} \\
\\ 
\texttt{**Punteggio:**} \\
\\ 
\texttt{Più vicini sono gli oggetti identici in entrambi i tavoli da gioco, più punti ottieni. Le penalità riducono i tuoi punti. Riesci a battere il record?} \\
\\ 
\\ 
\texttt{Per favore invia un messaggio all'altro giocatore per iniziare il gioco!} \\
            }
        }
    }
    & \\ \\

    \theutterance \stepcounter{utterance}  
    & \multicolumn{4}{p{0.6\linewidth}}{
        \cellcolor[rgb]{0.9,0.9,0.9}{
            \makecell[{{p{\linewidth}}}]{
                \texttt{\tiny{[P1$\rangle$GM]}}
                \texttt{COMUNICA: Ciao! Sono pronto per iniziare il gioco. Propongo di iniziare con un oggetto facile come "S" e provare a posizionarlo nella posizione (1,1) sul mio tavolo. Come vorresti procedere?} \\
            }
        }
    }
    & & \\ \\

    \theutterance \stepcounter{utterance}  
    & & & \multicolumn{2}{p{0.3\linewidth}}{
        \cellcolor[rgb]{0.9,0.9,0.9}{
            \makecell[{{p{\linewidth}}}]{
                \texttt{\tiny{[GM$|$GM]}}
                \texttt{COMUNICA: Ciao! Sono pronto per iniziare il gioco. Propongo di iniziare con un oggetto facile come "S" e provare a posizionarlo nella posizione (1,1) sul mio tavolo. Come vorresti procedere?} \\
            }
        }
    }
    & & \\ \\

    \theutterance \stepcounter{utterance}  
    & & \multicolumn{4}{p{0.6\linewidth}}{
        \cellcolor[rgb]{0.9,0.9,0.9}{
            \makecell[{{p{\linewidth}}}]{
                \texttt{\tiny{[GM$\rangle$P2]}}
                \texttt{Sono il tuo game master, e stai giocando un gioco collaborativo con la seguente griglia come tavolo da gioco:} \\
\\ 
\texttt{```} \\
\texttt{ 1234567} \\
\texttt{╔═══╤═╤═╗} \\
\texttt{║P◌A│◌│◌║ 1} \\
\texttt{╟─┬─┼─┴─╢ 2} \\
\texttt{║◌│D│◌◌◌║ 3} \\
\texttt{║◌├─┤◌┌─╢ 4} \\
\texttt{╟─┘◌│◌│O║ 5} \\
\texttt{║◌◌◌└┬┴─╢ 6} \\
\texttt{║◌S◌T│◌M║ 7} \\
\texttt{╚════╧══╝} \\
\\ 
\texttt{```} \\
\\ 
\texttt{* Il bordo superiore mostra le coordinate x che aumentano verso destra, e il bordo destro le coordinate y che aumentano verso il basso.} \\
\texttt{* The following objects are randomly placed on your grid: 'P', 'O', 'T', 'S', 'D', 'A', 'M'.} \\
\\ 
\\ 
\texttt{L'altro giocatore vede una variante del tavolo da gioco, dove gli oggetti sono posizionati in diverse posizioni casuali. Non puoi vedere il tavolo dell'altro giocatore, e loro non possono vedere il tuo.} \\
\\ 
\\ 
\texttt{**Obiettivo:**} \\
\\ 
\texttt{Entrambi i giocatori devono spostare gli oggetti sui rispettivi sfondi in modo che oggetti identici finiscano alle stesse coordinate. Devi comunicare con l'altro giocatore per concordare una configurazione obiettivo comune.} \\
\\ 
\\ 
\texttt{**Regole:**} \\
\\ 
\texttt{* In ogni turno, puoi inviare esattamente uno dei seguenti due comandi:} \\
\texttt{ 1. `COMUNICA: <MESSAGE>`: per inviare un messaggio (tutto fino alla prossima interruzione di riga) all'altro giocatore. Lo inoltrerò al tuo partner.} \\
\texttt{2. `SPOSTA: <OBJECT>, (<X>, <Y>)`: per spostare un oggetto in una nuova posizione, dove <X> è la colonna e <Y> è la riga. Ti farò sapere se la tua mossa è andata a buon fine o no.} \\
\texttt{* Se non rispetti il formato, o invii diversi comandi contemporaneamente, devo penalizzarti.} \\
\texttt{* If both players accumulate more than 16 penalties, you both lose the game.} \\
\texttt{* È vitale che comunichi con l'altro giocatore riguardo al tuo stato obiettivo! L'*unico* modo in cui puoi trasmettere la tua strategia all'altro giocatore è usando il comando `COMUNICA: <MESSAGE>`!} \\
\\ 
\\ 
\texttt{**Spostamento Oggetti**} \\
\\ 
\texttt{* Puoi spostare oggetti solo in celle entro i limiti della griglia. La cella di destinazione deve essere vuota, cioè deve contenere solo il simbolo '◌'.} \\
\texttt{* Se provi a spostare un oggetto in un posto che non è vuoto, o provi a spostarlo fuori dalla griglia, devo penalizzarti. Hai un altro tentativo.} \\
\texttt{* Prima di fare una mossa, controlla due volte che il posto di destinazione sia vuoto, e non contenga alcuna lettera, cornice, o linea!} \\
\\ 
\\ 
\texttt{**Fine del Gioco**} \\
\\ 
\texttt{Se pensi di aver raggiunto l'obiettivo di allineare tutti gli oggetti, puoi chiedere all'altro giocatore di finire il gioco inviando COMUNICA: finito?. Se l'altro giocatore ti chiede di finire il gioco, e rispondi con COMUNICA: finito!, the game will end.} \\
\\ 
\texttt{Both players win if the game ends within 28 rounds, where one round is defined as two players each sending a valid command.} \\
\\ 
\texttt{**Punteggio:**} \\
\\ 
\texttt{Più vicini sono gli oggetti identici in entrambi i tavoli da gioco, più punti ottieni. Le penalità riducono i tuoi punti. Riesci a battere il record?} \\
\\ 
\texttt{L'altro giocatore ha iniziato il gioco inviando questo messaggio:} \\
\\ 
\texttt{"Ciao! Sono pronto per iniziare il gioco. Propongo di iniziare con un oggetto facile come "S" e provare a posizionarlo nella posizione (1,1) sul mio tavolo. Come vorresti procedere?"} \\
\\ 
\texttt{Qual è il tuo primo comando?} \\
            }
        }
    }
    & \\ \\

    \theutterance \stepcounter{utterance}  
    & & & \multicolumn{4}{p{0.6\linewidth}}{
        \cellcolor[rgb]{0.9,0.9,0.9}{
            \makecell[{{p{\linewidth}}}]{
                \texttt{\tiny{[GM$\langle$P2]}}
                \texttt{Fantastico! Sono entusiasta di iniziare il gioco!} \\
\\ 
\texttt{Prima di tutto, vorrei chiarire che la mia griglia attuale è:} \\
\\ 
\texttt{```} \\
\texttt{ 1234567} \\
\texttt{╔═══╤═╤═╗} \\
\texttt{║P◌A│◌│◌║ 1} \\
\texttt{╟─┬─┼─┴─╢ 2} \\
\texttt{║◌│D│◌◌◌║ 3} \\
\texttt{║◌├─┤◌┌─╢ 4} \\
\texttt{╟─┘◌│◌│O║ 5} \\
\texttt{║◌◌◌└┬┴─╢ 6} \\
\texttt{║◌S◌T│◌M║ 7} \\
\texttt{╚════╧══╝} \\
\texttt{```} \\
\\ 
\texttt{Come proponi di iniziare con l'oggetto "S", mi sembra una buona idea! Tuttavia, vorrei sapere se il tuo tavolo ha già l'oggetto "S" nella posizione (1,1) o se è una proposta per entrambi di spostarlo lì.} \\
\\ 
\texttt{Il mio primo comando sarà:} \\
\\ 
\texttt{COMUNICA: "Ciao! La tua idea di iniziare con l'oggetto 'S' mi sembra buona. Tuttavia, vorrei sapere se il tuo tavolo ha già l'oggetto 'S' nella posizione (1,1) o se è una proposta per entrambi di spostarlo lì. Inoltre, vorrei condividere la mia griglia attuale con te, in modo che possiamo sincronizzare le nostre mosse."} \\
            }
        }
    }
     \\ \\

    \theutterance \stepcounter{utterance}  
    & & & \multicolumn{2}{p{0.3\linewidth}}{
        \cellcolor[rgb]{0.9,0.9,0.9}{
            \makecell[{{p{\linewidth}}}]{
                \texttt{\tiny{[GM$|$GM]}}
                \texttt{Fantastico! Sono entusiasta di iniziare il gioco!} \\
\\ 
\texttt{Prima di tutto, vorrei chiarire che la mia griglia attuale è:} \\
\\ 
\texttt{```} \\
\texttt{ 1234567} \\
\texttt{╔═══╤═╤═╗} \\
\texttt{║P◌A│◌│◌║ 1} \\
\texttt{╟─┬─┼─┴─╢ 2} \\
\texttt{║◌│D│◌◌◌║ 3} \\
\texttt{║◌├─┤◌┌─╢ 4} \\
\texttt{╟─┘◌│◌│O║ 5} \\
\texttt{║◌◌◌└┬┴─╢ 6} \\
\texttt{║◌S◌T│◌M║ 7} \\
\texttt{╚════╧══╝} \\
\texttt{```} \\
\\ 
\texttt{Come proponi di iniziare con l'oggetto "S", mi sembra una buona idea! Tuttavia, vorrei sapere se il tuo tavolo ha già l'oggetto "S" nella posizione (1,1) o se è una proposta per entrambi di spostarlo lì.} \\
\\ 
\texttt{Il mio primo comando sarà:} \\
\\ 
\texttt{COMUNICA: "Ciao! La tua idea di iniziare con l'oggetto 'S' mi sembra buona. Tuttavia, vorrei sapere se il tuo tavolo ha già l'oggetto 'S' nella posizione (1,1) o se è una proposta per entrambi di spostarlo lì. Inoltre, vorrei condividere la mia griglia attuale con te, in modo che possiamo sincronizzare le nostre mosse."} \\
            }
        }
    }
    & & \\ \\

    \theutterance \stepcounter{utterance}  
    & & & \multicolumn{2}{p{0.3\linewidth}}{
        \cellcolor[rgb]{0.9,0.9,0.9}{
            \makecell[{{p{\linewidth}}}]{
                \texttt{\tiny{[GM$|$GM]}}
                \texttt{Invalid format: head is not empty: 'Fantastico! Sono entusiasta di iniziare il gioco!} \\
\\ 
\texttt{Prima di tutto, vorrei chiarire che la mia griglia attuale è:} \\
\\ 
\\ 
\texttt{ 1234567} \\
\texttt{╔═══╤═╤═╗} \\
\texttt{║P◌A│◌│◌║ 1} \\
\texttt{╟─┬─┼─┴─╢ 2} \\
\texttt{║◌│D│◌◌◌║ 3} \\
\texttt{║◌├─┤◌┌─╢ 4} \\
\texttt{╟─┘◌│◌│O║ 5} \\
\texttt{║◌◌◌└┬┴─╢ 6} \\
\texttt{║◌S◌T│◌M║ 7} \\
\texttt{╚════╧══╝} \\
\\ 
\\ 
\texttt{Come proponi di iniziare con l'oggetto "S", mi sembra una buona idea! Tuttavia, vorrei sapere se il tuo tavolo ha già l'oggetto "S" nella posizione (1,1) o se è una proposta per entrambi di spostarlo lì.} \\
\\ 
\texttt{Il mio primo comando sarà:} \\
\\ 
\texttt{'} \\
            }
        }
    }
    & & \\ \\

    \theutterance \stepcounter{utterance}  
    & & \multicolumn{4}{p{0.6\linewidth}}{
        \cellcolor[rgb]{0.9,0.9,0.9}{
            \makecell[{{p{\linewidth}}}]{
                \texttt{\tiny{[GM$\rangle$P2]}}
                \texttt{Penalità: Il tuo messaggio non deve contenere nulla prima del comando!} \\
\\ 
\texttt{Assicurati che la tua risposta contenga solo `COMUNICA: <MESSAGE>` o `SPOSTA: <OBJECT>, (<X>, <Y>)`, e nient'altro!} \\
\\ 
\texttt{You have collectively accumulated 1 of 16 penalties.} \\
\texttt{Riprova!} \\
            }
        }
    }
    & \\ \\

    \theutterance \stepcounter{utterance}  
    & & & \multicolumn{4}{p{0.6\linewidth}}{
        \cellcolor[rgb]{0.9,0.9,0.9}{
            \makecell[{{p{\linewidth}}}]{
                \texttt{\tiny{[GM$\langle$P2]}}
                \texttt{COMUNICA: Ciao! La tua idea di iniziare con l'oggetto 'S' mi sembra buona. Tuttavia, vorrei sapere se il tuo tavolo ha già l'oggetto 'S' nella posizione (1,1) o se è una proposta per entrambi di spostarlo lì. Inoltre, vorrei condividere la mia griglia attuale con te, in modo che possiamo sincronizzare le nostre mosse.} \\
            }
        }
    }
     \\ \\

    \theutterance \stepcounter{utterance}  
    & & & \multicolumn{2}{p{0.3\linewidth}}{
        \cellcolor[rgb]{0.9,0.9,0.9}{
            \makecell[{{p{\linewidth}}}]{
                \texttt{\tiny{[GM$|$GM]}}
                \texttt{COMUNICA: Ciao! La tua idea di iniziare con l'oggetto 'S' mi sembra buona. Tuttavia, vorrei sapere se il tuo tavolo ha già l'oggetto 'S' nella posizione (1,1) o se è una proposta per entrambi di spostarlo lì. Inoltre, vorrei condividere la mia griglia attuale con te, in modo che possiamo sincronizzare le nostre mosse.} \\
            }
        }
    }
    & & \\ \\

    \theutterance \stepcounter{utterance}  
    & & \multicolumn{4}{p{0.6\linewidth}}{
        \cellcolor[rgb]{0.9,0.9,0.9}{
            \makecell[{{p{\linewidth}}}]{
                \texttt{\tiny{[P1$\langle$GM]}}
                \texttt{Il tuo messaggio è stato inoltrato all'altro giocatore.} \\
\\ 
\\ 
\texttt{You are currently playing round 1 of maximum 28.} \\
\texttt{You have collectively accumulated 1 of 16 penalties.} \\
\texttt{L'altro giocatore ha inviato questo messaggio:} \\
\\ 
\\ 
\texttt{"Ciao! La tua idea di iniziare con l'oggetto 'S' mi sembra buona. Tuttavia, vorrei sapere se il tuo tavolo ha già l'oggetto 'S' nella posizione (1,1) o se è una proposta per entrambi di spostarlo lì. Inoltre, vorrei condividere la mia griglia attuale con te, in modo che possiamo sincronizzare le nostre mosse."} \\
\\ 
\\ 
\texttt{Qual è il tuo prossimo comando?} \\
            }
        }
    }
    & \\ \\

    \theutterance \stepcounter{utterance}  
    & \multicolumn{4}{p{0.6\linewidth}}{
        \cellcolor[rgb]{0.9,0.9,0.9}{
            \makecell[{{p{\linewidth}}}]{
                \texttt{\tiny{[P1$\rangle$GM]}}
                \texttt{COMUNICA: Grazie per la risposta! Sì, l'oggetto 'S' si trova già nella posizione (1,1) sul mio tavolo, quindi non devo spostarlo. Vorrei vedere la tua griglia attuale per capire meglio la situazione e pianificare le prossime mosse. Per favore, condividi la tua griglia o descrivila in modo che io possa capire la disposizione degli oggetti. Inoltre, propongo di concentrarci su un altro oggetto dopo aver visto la tua griglia.} \\
            }
        }
    }
    & & \\ \\

    \theutterance \stepcounter{utterance}  
    & & & \multicolumn{2}{p{0.3\linewidth}}{
        \cellcolor[rgb]{0.9,0.9,0.9}{
            \makecell[{{p{\linewidth}}}]{
                \texttt{\tiny{[GM$|$GM]}}
                \texttt{COMUNICA: Grazie per la risposta! Sì, l'oggetto 'S' si trova già nella posizione (1,1) sul mio tavolo, quindi non devo spostarlo. Vorrei vedere la tua griglia attuale per capire meglio la situazione e pianificare le prossime mosse. Per favore, condividi la tua griglia o descrivila in modo che io possa capire la disposizione degli oggetti. Inoltre, propongo di concentrarci su un altro oggetto dopo aver visto la tua griglia.} \\
            }
        }
    }
    & & \\ \\

    \theutterance \stepcounter{utterance}  
    & & \multicolumn{4}{p{0.6\linewidth}}{
        \cellcolor[rgb]{0.9,0.9,0.9}{
            \makecell[{{p{\linewidth}}}]{
                \texttt{\tiny{[GM$\rangle$P2]}}
                \texttt{Il tuo messaggio è stato inoltrato all'altro giocatore.} \\
\\ 
\\ 
\texttt{You are currently playing round 2 of maximum 28.} \\
\texttt{You have collectively accumulated 1 of 16 penalties.} \\
\texttt{L'altro giocatore ha inviato questo messaggio:} \\
\\ 
\\ 
\texttt{"Grazie per la risposta! Sì, l'oggetto 'S' si trova già nella posizione (1,1) sul mio tavolo, quindi non devo spostarlo. Vorrei vedere la tua griglia attuale per capire meglio la situazione e pianificare le prossime mosse. Per favore, condividi la tua griglia o descrivila in modo che io possa capire la disposizione degli oggetti. Inoltre, propongo di concentrarci su un altro oggetto dopo aver visto la tua griglia."} \\
\\ 
\\ 
\texttt{Qual è il tuo prossimo comando?} \\
            }
        }
    }
    & \\ \\

    \theutterance \stepcounter{utterance}  
    & & & \multicolumn{4}{p{0.6\linewidth}}{
        \cellcolor[rgb]{0.9,0.9,0.9}{
            \makecell[{{p{\linewidth}}}]{
                \texttt{\tiny{[GM$\langle$P2]}}
                \texttt{COMUNICA: La mia griglia attuale è:} \\
\texttt{1234567} \\
\texttt{╔═══╤═╤═╗} \\
\texttt{║P◌A│◌│◌║ 1} \\
\texttt{╟─┬─┼─┴─╢ 2} \\
\texttt{║◌│D│◌◌◌║ 3} \\
\texttt{║◌├─┤◌┌─╢ 4} \\
\texttt{╟─┘◌│◌│O║ 5} \\
\texttt{║◌◌◌└┬┴─╢ 6} \\
\texttt{║◌S◌T│◌M║ 7} \\
\texttt{╚════╧══╝} \\
\texttt{Sono d'accordo di concentrarci su un altro oggetto dopo che hai visto la mia griglia. Propongo di considerare l'oggetto 'P' come prossimo obiettivo. Cosa ne pensi?} \\
            }
        }
    }
     \\ \\

    \theutterance \stepcounter{utterance}  
    & & & \multicolumn{2}{p{0.3\linewidth}}{
        \cellcolor[rgb]{0.9,0.9,0.9}{
            \makecell[{{p{\linewidth}}}]{
                \texttt{\tiny{[GM$|$GM]}}
                \texttt{COMUNICA: La mia griglia attuale è:} \\
\texttt{1234567} \\
\texttt{╔═══╤═╤═╗} \\
\texttt{║P◌A│◌│◌║ 1} \\
\texttt{╟─┬─┼─┴─╢ 2} \\
\texttt{║◌│D│◌◌◌║ 3} \\
\texttt{║◌├─┤◌┌─╢ 4} \\
\texttt{╟─┘◌│◌│O║ 5} \\
\texttt{║◌◌◌└┬┴─╢ 6} \\
\texttt{║◌S◌T│◌M║ 7} \\
\texttt{╚════╧══╝} \\
\texttt{Sono d'accordo di concentrarci su un altro oggetto dopo che hai visto la mia griglia. Propongo di considerare l'oggetto 'P' come prossimo obiettivo. Cosa ne pensi?} \\
            }
        }
    }
    & & \\ \\

    \theutterance \stepcounter{utterance}  
    & & & \multicolumn{2}{p{0.3\linewidth}}{
        \cellcolor[rgb]{0.9,0.9,0.9}{
            \makecell[{{p{\linewidth}}}]{
                \texttt{\tiny{[GM$|$GM]}}
                \texttt{Invalid format: tail is not empty: '} \\
\texttt{1234567} \\
\texttt{╔═══╤═╤═╗} \\
\texttt{║P◌A│◌│◌║ 1} \\
\texttt{╟─┬─┼─┴─╢ 2} \\
\texttt{║◌│D│◌◌◌║ 3} \\
\texttt{║◌├─┤◌┌─╢ 4} \\
\texttt{╟─┘◌│◌│O║ 5} \\
\texttt{║◌◌◌└┬┴─╢ 6} \\
\texttt{║◌S◌T│◌M║ 7} \\
\texttt{╚════╧══╝} \\
\texttt{Sono d'accordo di concentrarci su un altro oggetto dopo che hai visto la mia griglia. Propongo di considerare l'oggetto 'P' come prossimo obiettivo. Cosa ne pensi?} \\
            }
        }
    }
    & & \\ \\

    \theutterance \stepcounter{utterance}  
    & & \multicolumn{4}{p{0.6\linewidth}}{
        \cellcolor[rgb]{0.9,0.9,0.9}{
            \makecell[{{p{\linewidth}}}]{
                \texttt{\tiny{[GM$\rangle$P2]}}
                \texttt{Penalità: Il tuo messaggio non deve contenere nulla dopo il comando!} \\
\\ 
\texttt{Assicurati che la tua risposta contenga solo `COMUNICA: <MESSAGE>` o `SPOSTA: <OBJECT>, (<X>, <Y>)`, e nient'altro!} \\
\\ 
\texttt{You have collectively accumulated 2 of 16 penalties.} \\
\texttt{Riprova!} \\
            }
        }
    }
    & \\ \\

    \theutterance \stepcounter{utterance}  
    & & & \multicolumn{4}{p{0.6\linewidth}}{
        \cellcolor[rgb]{0.9,0.9,0.9}{
            \makecell[{{p{\linewidth}}}]{
                \texttt{\tiny{[GM$\langle$P2]}}
                \texttt{COMUNICA: La mia griglia attuale è: } \\
\texttt{1234567} \\
\texttt{╔═══╤═╤═╗} \\
\texttt{║P◌A│◌│◌║ 1} \\
\texttt{╟─┬─┼─┴─╢ 2} \\
\texttt{║◌│D│◌◌◌║ 3} \\
\texttt{║◌├─┤◌┌─╢ 4} \\
\texttt{╟─┘◌│◌│O║ 5} \\
\texttt{║◌◌◌└┬┴─╢ 6} \\
\texttt{║◌S◌T│◌M║ 7} \\
\texttt{╚════╧══╝} \\
\texttt{Sono d'accordo di concentrarci su un altro oggetto dopo che hai visto la mia griglia. Propongo di considerare l'oggetto 'P' come prossimo obiettivo. Cosa ne pensi?} \\
            }
        }
    }
     \\ \\

    \theutterance \stepcounter{utterance}  
    & & & \multicolumn{2}{p{0.3\linewidth}}{
        \cellcolor[rgb]{0.9,0.9,0.9}{
            \makecell[{{p{\linewidth}}}]{
                \texttt{\tiny{[GM$|$GM]}}
                \texttt{COMUNICA: La mia griglia attuale è: } \\
\texttt{1234567} \\
\texttt{╔═══╤═╤═╗} \\
\texttt{║P◌A│◌│◌║ 1} \\
\texttt{╟─┬─┼─┴─╢ 2} \\
\texttt{║◌│D│◌◌◌║ 3} \\
\texttt{║◌├─┤◌┌─╢ 4} \\
\texttt{╟─┘◌│◌│O║ 5} \\
\texttt{║◌◌◌└┬┴─╢ 6} \\
\texttt{║◌S◌T│◌M║ 7} \\
\texttt{╚════╧══╝} \\
\texttt{Sono d'accordo di concentrarci su un altro oggetto dopo che hai visto la mia griglia. Propongo di considerare l'oggetto 'P' come prossimo obiettivo. Cosa ne pensi?} \\
            }
        }
    }
    & & \\ \\

    \theutterance \stepcounter{utterance}  
    & & & \multicolumn{2}{p{0.3\linewidth}}{
        \cellcolor[rgb]{0.9,0.9,0.9}{
            \makecell[{{p{\linewidth}}}]{
                \texttt{\tiny{[GM$|$GM]}}
                \texttt{Invalid format: tail is not empty: '} \\
\texttt{1234567} \\
\texttt{╔═══╤═╤═╗} \\
\texttt{║P◌A│◌│◌║ 1} \\
\texttt{╟─┬─┼─┴─╢ 2} \\
\texttt{║◌│D│◌◌◌║ 3} \\
\texttt{║◌├─┤◌┌─╢ 4} \\
\texttt{╟─┘◌│◌│O║ 5} \\
\texttt{║◌◌◌└┬┴─╢ 6} \\
\texttt{║◌S◌T│◌M║ 7} \\
\texttt{╚════╧══╝} \\
\texttt{Sono d'accordo di concentrarci su un altro oggetto dopo che hai visto la mia griglia. Propongo di considerare l'oggetto 'P' come prossimo obiettivo. Cosa ne pensi?} \\
            }
        }
    }
    & & \\ \\

    \theutterance \stepcounter{utterance}  
    & & \multicolumn{4}{p{0.6\linewidth}}{
        \cellcolor[rgb]{0.9,0.9,0.9}{
            \makecell[{{p{\linewidth}}}]{
                \texttt{\tiny{[GM$\rangle$P2]}}
                \texttt{Penalità: Il tuo messaggio non deve contenere nulla dopo il comando!} \\
\\ 
\texttt{Assicurati che la tua risposta contenga solo `COMUNICA: <MESSAGE>` o `SPOSTA: <OBJECT>, (<X>, <Y>)`, e nient'altro!} \\
\\ 
\texttt{You have collectively accumulated 3 of 16 penalties.} \\
\texttt{Riprova!} \\
            }
        }
    }
    & \\ \\

    \theutterance \stepcounter{utterance}  
    & & & \multicolumn{4}{p{0.6\linewidth}}{
        \cellcolor[rgb]{0.9,0.9,0.9}{
            \makecell[{{p{\linewidth}}}]{
                \texttt{\tiny{[GM$\langle$P2]}}
                \texttt{COMUNICA: La mia griglia attuale è: } \\
\texttt{1234567} \\
\texttt{╔═══╤═╤═╗} \\
\texttt{║P◌A│◌│◌║ 1} \\
\texttt{╟─┬─┼─┴─╢ 2} \\
\texttt{║◌│D│◌◌◌║ 3} \\
\texttt{║◌├─┤◌┌─╢ 4} \\
\texttt{╟─┘◌│◌│O║ 5} \\
\texttt{║◌◌◌└┬┴─╢ 6} \\
\texttt{║◌S◌T│◌M║ 7} \\
\texttt{╚════╧══╝ Sono d'accordo di concentrarci su un altro oggetto dopo che hai visto la mia griglia. Propongo di considerare l'oggetto 'P' come prossimo obiettivo.} \\
            }
        }
    }
     \\ \\

    \theutterance \stepcounter{utterance}  
    & & & \multicolumn{2}{p{0.3\linewidth}}{
        \cellcolor[rgb]{0.9,0.9,0.9}{
            \makecell[{{p{\linewidth}}}]{
                \texttt{\tiny{[GM$|$GM]}}
                \texttt{COMUNICA: La mia griglia attuale è: } \\
\texttt{1234567} \\
\texttt{╔═══╤═╤═╗} \\
\texttt{║P◌A│◌│◌║ 1} \\
\texttt{╟─┬─┼─┴─╢ 2} \\
\texttt{║◌│D│◌◌◌║ 3} \\
\texttt{║◌├─┤◌┌─╢ 4} \\
\texttt{╟─┘◌│◌│O║ 5} \\
\texttt{║◌◌◌└┬┴─╢ 6} \\
\texttt{║◌S◌T│◌M║ 7} \\
\texttt{╚════╧══╝ Sono d'accordo di concentrarci su un altro oggetto dopo che hai visto la mia griglia. Propongo di considerare l'oggetto 'P' come prossimo obiettivo.} \\
            }
        }
    }
    & & \\ \\

    \theutterance \stepcounter{utterance}  
    & & & \multicolumn{2}{p{0.3\linewidth}}{
        \cellcolor[rgb]{0.9,0.9,0.9}{
            \makecell[{{p{\linewidth}}}]{
                \texttt{\tiny{[GM$|$GM]}}
                \texttt{Invalid format: tail is not empty: '} \\
\texttt{1234567} \\
\texttt{╔═══╤═╤═╗} \\
\texttt{║P◌A│◌│◌║ 1} \\
\texttt{╟─┬─┼─┴─╢ 2} \\
\texttt{║◌│D│◌◌◌║ 3} \\
\texttt{║◌├─┤◌┌─╢ 4} \\
\texttt{╟─┘◌│◌│O║ 5} \\
\texttt{║◌◌◌└┬┴─╢ 6} \\
\texttt{║◌S◌T│◌M║ 7} \\
\texttt{╚════╧══╝ Sono d'accordo di concentrarci su un altro oggetto dopo che hai visto la mia griglia. Propongo di considerare l'oggetto 'P' come prossimo obiettivo.} \\
            }
        }
    }
    & & \\ \\

    \theutterance \stepcounter{utterance}  
    & & \multicolumn{4}{p{0.6\linewidth}}{
        \cellcolor[rgb]{0.9,0.9,0.9}{
            \makecell[{{p{\linewidth}}}]{
                \texttt{\tiny{[GM$\rangle$P2]}}
                \texttt{Penalità: Il tuo messaggio non deve contenere nulla dopo il comando!} \\
\\ 
\texttt{Assicurati che la tua risposta contenga solo `COMUNICA: <MESSAGE>` o `SPOSTA: <OBJECT>, (<X>, <Y>)`, e nient'altro!} \\
\\ 
\texttt{You have collectively accumulated 4 of 16 penalties.} \\
\texttt{Riprova!} \\
            }
        }
    }
    & \\ \\

    \theutterance \stepcounter{utterance}  
    & & & \multicolumn{4}{p{0.6\linewidth}}{
        \cellcolor[rgb]{0.9,0.9,0.9}{
            \makecell[{{p{\linewidth}}}]{
                \texttt{\tiny{[GM$\langle$P2]}}
                \texttt{COMUNICA: La mia griglia attuale è: } \\
\texttt{1234567} \\
\texttt{╔═══╤═╤═╗} \\
\texttt{║P◌A│◌│◌║ 1} \\
\texttt{╟─┬─┼─┴─╢ 2} \\
\texttt{║◌│D│◌◌◌║ 3} \\
\texttt{║◌├─┤◌┌─╢ 4} \\
\texttt{╟─┘◌│◌│O║ 5} \\
\texttt{║◌◌◌└┬┴─╢ 6} \\
\texttt{║◌S◌T│◌M║ 7} \\
\texttt{╚════╧══╝ Sono d'accordo di concentrarci su un altro oggetto dopo che hai visto la mia griglia. Propongo di considerare l'oggetto 'P' come prossimo obiettivo. Cosa ne pensi?} \\
            }
        }
    }
     \\ \\

    \theutterance \stepcounter{utterance}  
    & & & \multicolumn{2}{p{0.3\linewidth}}{
        \cellcolor[rgb]{0.9,0.9,0.9}{
            \makecell[{{p{\linewidth}}}]{
                \texttt{\tiny{[GM$|$GM]}}
                \texttt{COMUNICA: La mia griglia attuale è: } \\
\texttt{1234567} \\
\texttt{╔═══╤═╤═╗} \\
\texttt{║P◌A│◌│◌║ 1} \\
\texttt{╟─┬─┼─┴─╢ 2} \\
\texttt{║◌│D│◌◌◌║ 3} \\
\texttt{║◌├─┤◌┌─╢ 4} \\
\texttt{╟─┘◌│◌│O║ 5} \\
\texttt{║◌◌◌└┬┴─╢ 6} \\
\texttt{║◌S◌T│◌M║ 7} \\
\texttt{╚════╧══╝ Sono d'accordo di concentrarci su un altro oggetto dopo che hai visto la mia griglia. Propongo di considerare l'oggetto 'P' come prossimo obiettivo. Cosa ne pensi?} \\
            }
        }
    }
    & & \\ \\

    \theutterance \stepcounter{utterance}  
    & & & \multicolumn{2}{p{0.3\linewidth}}{
        \cellcolor[rgb]{0.9,0.9,0.9}{
            \makecell[{{p{\linewidth}}}]{
                \texttt{\tiny{[GM$|$GM]}}
                \texttt{Invalid format: tail is not empty: '} \\
\texttt{1234567} \\
\texttt{╔═══╤═╤═╗} \\
\texttt{║P◌A│◌│◌║ 1} \\
\texttt{╟─┬─┼─┴─╢ 2} \\
\texttt{║◌│D│◌◌◌║ 3} \\
\texttt{║◌├─┤◌┌─╢ 4} \\
\texttt{╟─┘◌│◌│O║ 5} \\
\texttt{║◌◌◌└┬┴─╢ 6} \\
\texttt{║◌S◌T│◌M║ 7} \\
\texttt{╚════╧══╝ Sono d'accordo di concentrarci su un altro oggetto dopo che hai visto la mia griglia. Propongo di considerare l'oggetto 'P' come prossimo obiettivo. Cosa ne pensi?} \\
            }
        }
    }
    & & \\ \\

    \theutterance \stepcounter{utterance}  
    & & \multicolumn{4}{p{0.6\linewidth}}{
        \cellcolor[rgb]{0.9,0.9,0.9}{
            \makecell[{{p{\linewidth}}}]{
                \texttt{\tiny{[GM$\rangle$P2]}}
                \texttt{Penalità: Il tuo messaggio non deve contenere nulla dopo il comando!} \\
\\ 
\texttt{Assicurati che la tua risposta contenga solo `COMUNICA: <MESSAGE>` o `SPOSTA: <OBJECT>, (<X>, <Y>)`, e nient'altro!} \\
\\ 
\texttt{You have collectively accumulated 5 of 16 penalties.} \\
\texttt{Riprova!} \\
            }
        }
    }
    & \\ \\

    \theutterance \stepcounter{utterance}  
    & & & \multicolumn{4}{p{0.6\linewidth}}{
        \cellcolor[rgb]{0.9,0.9,0.9}{
            \makecell[{{p{\linewidth}}}]{
                \texttt{\tiny{[GM$\langle$P2]}}
                \texttt{COMUNICA: La mia griglia attuale è: } \\
\texttt{1234567} \\
\texttt{╔═══╤═╤═╗} \\
\texttt{║P◌A│◌│◌║ 1} \\
\texttt{╟─┬─┼─┴─╢ 2} \\
\texttt{║◌│D│◌◌◌║ 3} \\
\texttt{║◌├─┤◌┌─╢ 4} \\
\texttt{╟─┘◌│◌│O║ 5} \\
\texttt{║◌◌◌└┬┴─╢ 6} \\
\texttt{║◌S◌T│◌M║ 7} \\
\texttt{╚════╧══╝ Sono d'accordo di concentrarci su un altro oggetto dopo che hai visto la mia griglia. Propongo di considerare l'oggetto 'P' come prossimo obiettivo.} \\
            }
        }
    }
     \\ \\

    \theutterance \stepcounter{utterance}  
    & & & \multicolumn{2}{p{0.3\linewidth}}{
        \cellcolor[rgb]{0.9,0.9,0.9}{
            \makecell[{{p{\linewidth}}}]{
                \texttt{\tiny{[GM$|$GM]}}
                \texttt{COMUNICA: La mia griglia attuale è: } \\
\texttt{1234567} \\
\texttt{╔═══╤═╤═╗} \\
\texttt{║P◌A│◌│◌║ 1} \\
\texttt{╟─┬─┼─┴─╢ 2} \\
\texttt{║◌│D│◌◌◌║ 3} \\
\texttt{║◌├─┤◌┌─╢ 4} \\
\texttt{╟─┘◌│◌│O║ 5} \\
\texttt{║◌◌◌└┬┴─╢ 6} \\
\texttt{║◌S◌T│◌M║ 7} \\
\texttt{╚════╧══╝ Sono d'accordo di concentrarci su un altro oggetto dopo che hai visto la mia griglia. Propongo di considerare l'oggetto 'P' come prossimo obiettivo.} \\
            }
        }
    }
    & & \\ \\

    \theutterance \stepcounter{utterance}  
    & & & \multicolumn{2}{p{0.3\linewidth}}{
        \cellcolor[rgb]{0.9,0.9,0.9}{
            \makecell[{{p{\linewidth}}}]{
                \texttt{\tiny{[GM$|$GM]}}
                \texttt{Invalid format: tail is not empty: '} \\
\texttt{1234567} \\
\texttt{╔═══╤═╤═╗} \\
\texttt{║P◌A│◌│◌║ 1} \\
\texttt{╟─┬─┼─┴─╢ 2} \\
\texttt{║◌│D│◌◌◌║ 3} \\
\texttt{║◌├─┤◌┌─╢ 4} \\
\texttt{╟─┘◌│◌│O║ 5} \\
\texttt{║◌◌◌└┬┴─╢ 6} \\
\texttt{║◌S◌T│◌M║ 7} \\
\texttt{╚════╧══╝ Sono d'accordo di concentrarci su un altro oggetto dopo che hai visto la mia griglia. Propongo di considerare l'oggetto 'P' come prossimo obiettivo.} \\
            }
        }
    }
    & & \\ \\

    \theutterance \stepcounter{utterance}  
    & & \multicolumn{4}{p{0.6\linewidth}}{
        \cellcolor[rgb]{0.9,0.9,0.9}{
            \makecell[{{p{\linewidth}}}]{
                \texttt{\tiny{[GM$\rangle$P2]}}
                \texttt{Penalità: Il tuo messaggio non deve contenere nulla dopo il comando!} \\
\\ 
\texttt{Assicurati che la tua risposta contenga solo `COMUNICA: <MESSAGE>` o `SPOSTA: <OBJECT>, (<X>, <Y>)`, e nient'altro!} \\
\\ 
\texttt{You have collectively accumulated 6 of 16 penalties.} \\
\texttt{Riprova!} \\
            }
        }
    }
    & \\ \\

    \theutterance \stepcounter{utterance}  
    & & & \multicolumn{4}{p{0.6\linewidth}}{
        \cellcolor[rgb]{0.9,0.9,0.9}{
            \makecell[{{p{\linewidth}}}]{
                \texttt{\tiny{[GM$\langle$P2]}}
                \texttt{COMUNICA: La mia griglia attuale è: } \\
\texttt{1234567} \\
\texttt{╔═══╤═╤═╗} \\
\texttt{║P◌A│◌│◌║ 1} \\
\texttt{╟─┬─┼─┴─╢ 2} \\
\texttt{║◌│D│◌◌◌║ 3} \\
\texttt{║◌├─┤◌┌─╢ 4} \\
\texttt{╟─┘◌│◌│O║ 5} \\
\texttt{║◌◌◌└┬┴─╢ 6} \\
\texttt{║◌S◌T│◌M║ 7} \\
\texttt{╚════╧══╝} \\
            }
        }
    }
     \\ \\

    \theutterance \stepcounter{utterance}  
    & & & \multicolumn{2}{p{0.3\linewidth}}{
        \cellcolor[rgb]{0.9,0.9,0.9}{
            \makecell[{{p{\linewidth}}}]{
                \texttt{\tiny{[GM$|$GM]}}
                \texttt{COMUNICA: La mia griglia attuale è: } \\
\texttt{1234567} \\
\texttt{╔═══╤═╤═╗} \\
\texttt{║P◌A│◌│◌║ 1} \\
\texttt{╟─┬─┼─┴─╢ 2} \\
\texttt{║◌│D│◌◌◌║ 3} \\
\texttt{║◌├─┤◌┌─╢ 4} \\
\texttt{╟─┘◌│◌│O║ 5} \\
\texttt{║◌◌◌└┬┴─╢ 6} \\
\texttt{║◌S◌T│◌M║ 7} \\
\texttt{╚════╧══╝} \\
            }
        }
    }
    & & \\ \\

    \theutterance \stepcounter{utterance}  
    & & & \multicolumn{2}{p{0.3\linewidth}}{
        \cellcolor[rgb]{0.9,0.9,0.9}{
            \makecell[{{p{\linewidth}}}]{
                \texttt{\tiny{[GM$|$GM]}}
                \texttt{Invalid format: tail is not empty: '} \\
\texttt{1234567} \\
\texttt{╔═══╤═╤═╗} \\
\texttt{║P◌A│◌│◌║ 1} \\
\texttt{╟─┬─┼─┴─╢ 2} \\
\texttt{║◌│D│◌◌◌║ 3} \\
\texttt{║◌├─┤◌┌─╢ 4} \\
\texttt{╟─┘◌│◌│O║ 5} \\
\texttt{║◌◌◌└┬┴─╢ 6} \\
\texttt{║◌S◌T│◌M║ 7} \\
\texttt{╚════╧══╝} \\
            }
        }
    }
    & & \\ \\

    \theutterance \stepcounter{utterance}  
    & & \multicolumn{4}{p{0.6\linewidth}}{
        \cellcolor[rgb]{0.9,0.9,0.9}{
            \makecell[{{p{\linewidth}}}]{
                \texttt{\tiny{[GM$\rangle$P2]}}
                \texttt{Penalità: Il tuo messaggio non deve contenere nulla dopo il comando!} \\
\\ 
\texttt{Assicurati che la tua risposta contenga solo `COMUNICA: <MESSAGE>` o `SPOSTA: <OBJECT>, (<X>, <Y>)`, e nient'altro!} \\
\\ 
\texttt{You have collectively accumulated 7 of 16 penalties.} \\
\texttt{Riprova!} \\
            }
        }
    }
    & \\ \\

    \theutterance \stepcounter{utterance}  
    & & & \multicolumn{4}{p{0.6\linewidth}}{
        \cellcolor[rgb]{0.9,0.9,0.9}{
            \makecell[{{p{\linewidth}}}]{
                \texttt{\tiny{[GM$\langle$P2]}}
                \texttt{COMUNICA: La mia griglia attuale è: 1234567 ╔═══╤═╤═╗ ║P◌A│◌│◌║ 1 ╟─┬─┼─┴─╢ 2 ║◌│D│◌◌◌║ 3 ║◌├─┤◌┌─╢ 4 ╟─┘◌│◌│O║ 5 ║◌◌◌└┬┴─╢ 6 ║◌S◌T│◌M║ 7 ╚════╧══╝} \\
            }
        }
    }
     \\ \\

    \theutterance \stepcounter{utterance}  
    & & & \multicolumn{2}{p{0.3\linewidth}}{
        \cellcolor[rgb]{0.9,0.9,0.9}{
            \makecell[{{p{\linewidth}}}]{
                \texttt{\tiny{[GM$|$GM]}}
                \texttt{COMUNICA: La mia griglia attuale è: 1234567 ╔═══╤═╤═╗ ║P◌A│◌│◌║ 1 ╟─┬─┼─┴─╢ 2 ║◌│D│◌◌◌║ 3 ║◌├─┤◌┌─╢ 4 ╟─┘◌│◌│O║ 5 ║◌◌◌└┬┴─╢ 6 ║◌S◌T│◌M║ 7 ╚════╧══╝} \\
            }
        }
    }
    & & \\ \\

    \theutterance \stepcounter{utterance}  
    & & \multicolumn{4}{p{0.6\linewidth}}{
        \cellcolor[rgb]{0.9,0.9,0.9}{
            \makecell[{{p{\linewidth}}}]{
                \texttt{\tiny{[P1$\langle$GM]}}
                \texttt{Il tuo messaggio è stato inoltrato all'altro giocatore.} \\
\\ 
\\ 
\texttt{You are currently playing round 2 of maximum 28.} \\
\texttt{You have collectively accumulated 7 of 16 penalties.} \\
\texttt{L'altro giocatore ha inviato questo messaggio:} \\
\\ 
\\ 
\texttt{"La mia griglia attuale è: 1234567 ╔═══╤═╤═╗ ║P◌A│◌│◌║ 1 ╟─┬─┼─┴─╢ 2 ║◌│D│◌◌◌║ 3 ║◌├─┤◌┌─╢ 4 ╟─┘◌│◌│O║ 5 ║◌◌◌└┬┴─╢ 6 ║◌S◌T│◌M║ 7 ╚════╧══╝"} \\
\\ 
\\ 
\texttt{Qual è il tuo prossimo comando?} \\
            }
        }
    }
    & \\ \\

    \theutterance \stepcounter{utterance}  
    & \multicolumn{4}{p{0.6\linewidth}}{
        \cellcolor[rgb]{0.9,0.9,0.9}{
            \makecell[{{p{\linewidth}}}]{
                \texttt{\tiny{[P1$\rangle$GM]}}
                \texttt{COMUNICA: Grazie per aver condiviso la tua griglia! Ora che posso vedere la disposizione degli oggetti sul tuo tavolo, posso pianificare meglio le prossime mosse. Noterò che l'oggetto 'S' si trova nella mia griglia alla posizione (1,1) e nella tua griglia alla posizione (7,7). Questo mi suggerisce che potremmo dover spostare gli oggetti in modo che le posizioni siano simili. Vorrei proporre di spostare gli oggetti 'P' e 'D' come prossimo passo, poiché sembrano essere in posizioni diverse rispetto alla mia griglia. Cosa ne pensi? Dovremmo provare a spostare 'P' e 'D' in posizioni più simili tra le nostre griglie?} \\
            }
        }
    }
    & & \\ \\

    \theutterance \stepcounter{utterance}  
    & & & \multicolumn{2}{p{0.3\linewidth}}{
        \cellcolor[rgb]{0.9,0.9,0.9}{
            \makecell[{{p{\linewidth}}}]{
                \texttt{\tiny{[GM$|$GM]}}
                \texttt{COMUNICA: Grazie per aver condiviso la tua griglia! Ora che posso vedere la disposizione degli oggetti sul tuo tavolo, posso pianificare meglio le prossime mosse. Noterò che l'oggetto 'S' si trova nella mia griglia alla posizione (1,1) e nella tua griglia alla posizione (7,7). Questo mi suggerisce che potremmo dover spostare gli oggetti in modo che le posizioni siano simili. Vorrei proporre di spostare gli oggetti 'P' e 'D' come prossimo passo, poiché sembrano essere in posizioni diverse rispetto alla mia griglia. Cosa ne pensi? Dovremmo provare a spostare 'P' e 'D' in posizioni più simili tra le nostre griglie?} \\
            }
        }
    }
    & & \\ \\

    \theutterance \stepcounter{utterance}  
    & & \multicolumn{4}{p{0.6\linewidth}}{
        \cellcolor[rgb]{0.9,0.9,0.9}{
            \makecell[{{p{\linewidth}}}]{
                \texttt{\tiny{[GM$\rangle$P2]}}
                \texttt{Il tuo messaggio è stato inoltrato all'altro giocatore.} \\
\\ 
\\ 
\texttt{You are currently playing round 3 of maximum 28.} \\
\texttt{You have collectively accumulated 7 of 16 penalties.} \\
\texttt{L'altro giocatore ha inviato questo messaggio:} \\
\\ 
\\ 
\texttt{"Grazie per aver condiviso la tua griglia! Ora che posso vedere la disposizione degli oggetti sul tuo tavolo, posso pianificare meglio le prossime mosse. Noterò che l'oggetto 'S' si trova nella mia griglia alla posizione (1,1) e nella tua griglia alla posizione (7,7). Questo mi suggerisce che potremmo dover spostare gli oggetti in modo che le posizioni siano simili. Vorrei proporre di spostare gli oggetti 'P' e 'D' come prossimo passo, poiché sembrano essere in posizioni diverse rispetto alla mia griglia. Cosa ne pensi? Dovremmo provare a spostare 'P' e 'D' in posizioni più simili tra le nostre griglie?"} \\
\\ 
\\ 
\texttt{Qual è il tuo prossimo comando?} \\
            }
        }
    }
    & \\ \\

    \theutterance \stepcounter{utterance}  
    & & & \multicolumn{4}{p{0.6\linewidth}}{
        \cellcolor[rgb]{0.9,0.9,0.9}{
            \makecell[{{p{\linewidth}}}]{
                \texttt{\tiny{[GM$\langle$P2]}}
                \texttt{COMUNICA: Sono d'accordo che spostare gli oggetti 'P' e 'D' potrebbe essere un buon prossimo passo, poiché sono in posizioni diverse rispetto alla tua griglia. Tuttavia, vorrei assicurarmi di comprendere bene la tua proposta prima di effettuare qualsiasi mossa. Potresti per favore fornire maggiori dettagli su dove proponi di spostare 'P' e 'D' e come ciò aiuterà a sincronizzare le nostre griglie?} \\
            }
        }
    }
     \\ \\

    \theutterance \stepcounter{utterance}  
    & & & \multicolumn{2}{p{0.3\linewidth}}{
        \cellcolor[rgb]{0.9,0.9,0.9}{
            \makecell[{{p{\linewidth}}}]{
                \texttt{\tiny{[GM$|$GM]}}
                \texttt{COMUNICA: Sono d'accordo che spostare gli oggetti 'P' e 'D' potrebbe essere un buon prossimo passo, poiché sono in posizioni diverse rispetto alla tua griglia. Tuttavia, vorrei assicurarmi di comprendere bene la tua proposta prima di effettuare qualsiasi mossa. Potresti per favore fornire maggiori dettagli su dove proponi di spostare 'P' e 'D' e come ciò aiuterà a sincronizzare le nostre griglie?} \\
            }
        }
    }
    & & \\ \\

    \theutterance \stepcounter{utterance}  
    & & \multicolumn{4}{p{0.6\linewidth}}{
        \cellcolor[rgb]{0.9,0.9,0.9}{
            \makecell[{{p{\linewidth}}}]{
                \texttt{\tiny{[P1$\langle$GM]}}
                \texttt{Il tuo messaggio è stato inoltrato all'altro giocatore.} \\
\\ 
\\ 
\texttt{You are currently playing round 3 of maximum 28.} \\
\texttt{You have collectively accumulated 7 of 16 penalties.} \\
\texttt{L'altro giocatore ha inviato questo messaggio:} \\
\\ 
\\ 
\texttt{"Sono d'accordo che spostare gli oggetti 'P' e 'D' potrebbe essere un buon prossimo passo, poiché sono in posizioni diverse rispetto alla tua griglia. Tuttavia, vorrei assicurarmi di comprendere bene la tua proposta prima di effettuare qualsiasi mossa. Potresti per favore fornire maggiori dettagli su dove proponi di spostare 'P' e 'D' e come ciò aiuterà a sincronizzare le nostre griglie?"} \\
\\ 
\\ 
\texttt{Qual è il tuo prossimo comando?} \\
            }
        }
    }
    & \\ \\

    \theutterance \stepcounter{utterance}  
    & \multicolumn{4}{p{0.6\linewidth}}{
        \cellcolor[rgb]{0.9,0.9,0.9}{
            \makecell[{{p{\linewidth}}}]{
                \texttt{\tiny{[P1$\rangle$GM]}}
                \texttt{COMUNICA: Grazie per la tua risposta! Sono felice di fornire maggiori dettagli sulla mia proposta. Noterò che l'oggetto 'P' si trova nella mia griglia alla posizione (4,4) e nella tua griglia alla posizione (1,1). Vorrei proporre di spostare il 'P' nella tua griglia alla posizione (4,4) per allinearla con la mia griglia. Per quanto riguarda l'oggetto 'D', si trova nella mia griglia alla posizione (4,5) e nella tua griglia alla posizione (3,3). Potremmo provare a spostare il 'D' nella tua griglia alla posizione (4,5) per allinearla con la mia griglia. Tuttavia, prima di effettuare qualsiasi mossa, vorrei assicurarmi che le posizioni di destinazione siano vuote e che non ci siano conflitti con altri oggetti. Cosa ne pensi di questa proposta? Dovremmo provare a spostare 'P' e 'D' secondo la mia proposta o hai una strategia alternativa?} \\
            }
        }
    }
    & & \\ \\

    \theutterance \stepcounter{utterance}  
    & & & \multicolumn{2}{p{0.3\linewidth}}{
        \cellcolor[rgb]{0.9,0.9,0.9}{
            \makecell[{{p{\linewidth}}}]{
                \texttt{\tiny{[GM$|$GM]}}
                \texttt{COMUNICA: Grazie per la tua risposta! Sono felice di fornire maggiori dettagli sulla mia proposta. Noterò che l'oggetto 'P' si trova nella mia griglia alla posizione (4,4) e nella tua griglia alla posizione (1,1). Vorrei proporre di spostare il 'P' nella tua griglia alla posizione (4,4) per allinearla con la mia griglia. Per quanto riguarda l'oggetto 'D', si trova nella mia griglia alla posizione (4,5) e nella tua griglia alla posizione (3,3). Potremmo provare a spostare il 'D' nella tua griglia alla posizione (4,5) per allinearla con la mia griglia. Tuttavia, prima di effettuare qualsiasi mossa, vorrei assicurarmi che le posizioni di destinazione siano vuote e che non ci siano conflitti con altri oggetti. Cosa ne pensi di questa proposta? Dovremmo provare a spostare 'P' e 'D' secondo la mia proposta o hai una strategia alternativa?} \\
            }
        }
    }
    & & \\ \\

    \theutterance \stepcounter{utterance}  
    & & \multicolumn{4}{p{0.6\linewidth}}{
        \cellcolor[rgb]{0.9,0.9,0.9}{
            \makecell[{{p{\linewidth}}}]{
                \texttt{\tiny{[GM$\rangle$P2]}}
                \texttt{Il tuo messaggio è stato inoltrato all'altro giocatore.} \\
\\ 
\\ 
\texttt{You are currently playing round 4 of maximum 28.} \\
\texttt{You have collectively accumulated 7 of 16 penalties.} \\
\texttt{L'altro giocatore ha inviato questo messaggio:} \\
\\ 
\\ 
\texttt{"Grazie per la tua risposta! Sono felice di fornire maggiori dettagli sulla mia proposta. Noterò che l'oggetto 'P' si trova nella mia griglia alla posizione (4,4) e nella tua griglia alla posizione (1,1). Vorrei proporre di spostare il 'P' nella tua griglia alla posizione (4,4) per allinearla con la mia griglia. Per quanto riguarda l'oggetto 'D', si trova nella mia griglia alla posizione (4,5) e nella tua griglia alla posizione (3,3). Potremmo provare a spostare il 'D' nella tua griglia alla posizione (4,5) per allinearla con la mia griglia. Tuttavia, prima di effettuare qualsiasi mossa, vorrei assicurarmi che le posizioni di destinazione siano vuote e che non ci siano conflitti con altri oggetti. Cosa ne pensi di questa proposta? Dovremmo provare a spostare 'P' e 'D' secondo la mia proposta o hai una strategia alternativa?"} \\
\\ 
\\ 
\texttt{Qual è il tuo prossimo comando?} \\
            }
        }
    }
    & \\ \\

    \theutterance \stepcounter{utterance}  
    & & & \multicolumn{4}{p{0.6\linewidth}}{
        \cellcolor[rgb]{0.9,0.9,0.9}{
            \makecell[{{p{\linewidth}}}]{
                \texttt{\tiny{[GM$\langle$P2]}}
                \texttt{COMUNICA: La tua proposta di spostare 'P' alla posizione (4,4) e 'D' alla posizione (4,5) sulla mia griglia sembra ragionevole, poiché potrebbe aiutare ad allineare le nostre griglie. Tuttavia, devo controllare la disposizione degli oggetti sulla mia griglia per assicurarmi che le posizioni di destinazione siano vuote. Una volta che avrò verificato, potrò decidere se procedere con la tua proposta o suggerire una strategia alternativa.} \\
            }
        }
    }
     \\ \\

    \theutterance \stepcounter{utterance}  
    & & & \multicolumn{2}{p{0.3\linewidth}}{
        \cellcolor[rgb]{0.9,0.9,0.9}{
            \makecell[{{p{\linewidth}}}]{
                \texttt{\tiny{[GM$|$GM]}}
                \texttt{COMUNICA: La tua proposta di spostare 'P' alla posizione (4,4) e 'D' alla posizione (4,5) sulla mia griglia sembra ragionevole, poiché potrebbe aiutare ad allineare le nostre griglie. Tuttavia, devo controllare la disposizione degli oggetti sulla mia griglia per assicurarmi che le posizioni di destinazione siano vuote. Una volta che avrò verificato, potrò decidere se procedere con la tua proposta o suggerire una strategia alternativa.} \\
            }
        }
    }
    & & \\ \\

    \theutterance \stepcounter{utterance}  
    & & \multicolumn{4}{p{0.6\linewidth}}{
        \cellcolor[rgb]{0.9,0.9,0.9}{
            \makecell[{{p{\linewidth}}}]{
                \texttt{\tiny{[P1$\langle$GM]}}
                \texttt{Il tuo messaggio è stato inoltrato all'altro giocatore.} \\
\\ 
\\ 
\texttt{You are currently playing round 4 of maximum 28.} \\
\texttt{You have collectively accumulated 7 of 16 penalties.} \\
\texttt{L'altro giocatore ha inviato questo messaggio:} \\
\\ 
\\ 
\texttt{"La tua proposta di spostare 'P' alla posizione (4,4) e 'D' alla posizione (4,5) sulla mia griglia sembra ragionevole, poiché potrebbe aiutare ad allineare le nostre griglie. Tuttavia, devo controllare la disposizione degli oggetti sulla mia griglia per assicurarmi che le posizioni di destinazione siano vuote. Una volta che avrò verificato, potrò decidere se procedere con la tua proposta o suggerire una strategia alternativa."} \\
\\ 
\\ 
\texttt{Qual è il tuo prossimo comando?} \\
            }
        }
    }
    & \\ \\

    \theutterance \stepcounter{utterance}  
    & \multicolumn{4}{p{0.6\linewidth}}{
        \cellcolor[rgb]{0.9,0.9,0.9}{
            \makecell[{{p{\linewidth}}}]{
                \texttt{\tiny{[P1$\rangle$GM]}}
                \texttt{COMUNICA: Grazie per la tua risposta! Sono felice di aspettare che tu verifichi la disposizione degli oggetti sulla tua griglia. Nel frattempo, vorrei chiederti se puoi verificare anche la posizione degli oggetti 'T', 'O' e 'M' sulla tua griglia, poiché potrebbero essere importanti per la strategia generale. Inoltre, ti suggerisco di comunicarmi se ci sono altre posizioni che potrebbero essere utili per allineare le nostre griglie, in modo da poter pianificare le prossime mosse in modo più efficace. Quanto tempo pensi di impiegare per verificare la disposizione degli oggetti sulla tua griglia? Possiamo procedere con le mosse una volta che avrai finito?} \\
            }
        }
    }
    & & \\ \\

    \theutterance \stepcounter{utterance}  
    & & & \multicolumn{2}{p{0.3\linewidth}}{
        \cellcolor[rgb]{0.9,0.9,0.9}{
            \makecell[{{p{\linewidth}}}]{
                \texttt{\tiny{[GM$|$GM]}}
                \texttt{COMUNICA: Grazie per la tua risposta! Sono felice di aspettare che tu verifichi la disposizione degli oggetti sulla tua griglia. Nel frattempo, vorrei chiederti se puoi verificare anche la posizione degli oggetti 'T', 'O' e 'M' sulla tua griglia, poiché potrebbero essere importanti per la strategia generale. Inoltre, ti suggerisco di comunicarmi se ci sono altre posizioni che potrebbero essere utili per allineare le nostre griglie, in modo da poter pianificare le prossime mosse in modo più efficace. Quanto tempo pensi di impiegare per verificare la disposizione degli oggetti sulla tua griglia? Possiamo procedere con le mosse una volta che avrai finito?} \\
            }
        }
    }
    & & \\ \\

    \theutterance \stepcounter{utterance}  
    & & \multicolumn{4}{p{0.6\linewidth}}{
        \cellcolor[rgb]{0.9,0.9,0.9}{
            \makecell[{{p{\linewidth}}}]{
                \texttt{\tiny{[GM$\rangle$P2]}}
                \texttt{Il tuo messaggio è stato inoltrato all'altro giocatore.} \\
\\ 
\\ 
\texttt{You are currently playing round 5 of maximum 28.} \\
\texttt{You have collectively accumulated 7 of 16 penalties.} \\
\texttt{L'altro giocatore ha inviato questo messaggio:} \\
\\ 
\\ 
\texttt{"Grazie per la tua risposta! Sono felice di aspettare che tu verifichi la disposizione degli oggetti sulla tua griglia. Nel frattempo, vorrei chiederti se puoi verificare anche la posizione degli oggetti 'T', 'O' e 'M' sulla tua griglia, poiché potrebbero essere importanti per la strategia generale. Inoltre, ti suggerisco di comunicarmi se ci sono altre posizioni che potrebbero essere utili per allineare le nostre griglie, in modo da poter pianificare le prossime mosse in modo più efficace. Quanto tempo pensi di impiegare per verificare la disposizione degli oggetti sulla tua griglia? Possiamo procedere con le mosse una volta che avrai finito?"} \\
\\ 
\\ 
\texttt{Qual è il tuo prossimo comando?} \\
            }
        }
    }
    & \\ \\

    \theutterance \stepcounter{utterance}  
    & & & \multicolumn{4}{p{0.6\linewidth}}{
        \cellcolor[rgb]{0.9,0.9,0.9}{
            \makecell[{{p{\linewidth}}}]{
                \texttt{\tiny{[GM$\langle$P2]}}
                \texttt{COMUNICA: Ho verificato la disposizione degli oggetti sulla mia griglia e posso confermare che le posizioni proposte per 'P' e 'D' sono disponibili. La posizione di 'T' è (7,7), 'O' è (5,5) e 'M' è (7,7). Sono pronto a procedere con le mosse per allineare le nostre griglie. Spostare 'P' alla posizione (4,4) e 'D' alla posizione (4,5) sembra una buona strategia. Sono pronto a eseguire le mosse necessarie per allineare le griglie.} \\
            }
        }
    }
     \\ \\

    \theutterance \stepcounter{utterance}  
    & & & \multicolumn{2}{p{0.3\linewidth}}{
        \cellcolor[rgb]{0.9,0.9,0.9}{
            \makecell[{{p{\linewidth}}}]{
                \texttt{\tiny{[GM$|$GM]}}
                \texttt{COMUNICA: Ho verificato la disposizione degli oggetti sulla mia griglia e posso confermare che le posizioni proposte per 'P' e 'D' sono disponibili. La posizione di 'T' è (7,7), 'O' è (5,5) e 'M' è (7,7). Sono pronto a procedere con le mosse per allineare le nostre griglie. Spostare 'P' alla posizione (4,4) e 'D' alla posizione (4,5) sembra una buona strategia. Sono pronto a eseguire le mosse necessarie per allineare le griglie.} \\
            }
        }
    }
    & & \\ \\

    \theutterance \stepcounter{utterance}  
    & & \multicolumn{4}{p{0.6\linewidth}}{
        \cellcolor[rgb]{0.9,0.9,0.9}{
            \makecell[{{p{\linewidth}}}]{
                \texttt{\tiny{[P1$\langle$GM]}}
                \texttt{Il tuo messaggio è stato inoltrato all'altro giocatore.} \\
\\ 
\\ 
\texttt{You are currently playing round 5 of maximum 28.} \\
\texttt{You have collectively accumulated 7 of 16 penalties.} \\
\texttt{L'altro giocatore ha inviato questo messaggio:} \\
\\ 
\\ 
\texttt{"Ho verificato la disposizione degli oggetti sulla mia griglia e posso confermare che le posizioni proposte per 'P' e 'D' sono disponibili. La posizione di 'T' è (7,7), 'O' è (5,5) e 'M' è (7,7). Sono pronto a procedere con le mosse per allineare le nostre griglie. Spostare 'P' alla posizione (4,4) e 'D' alla posizione (4,5) sembra una buona strategia. Sono pronto a eseguire le mosse necessarie per allineare le griglie."} \\
\\ 
\\ 
\texttt{Qual è il tuo prossimo comando?} \\
            }
        }
    }
    & \\ \\

    \theutterance \stepcounter{utterance}  
    & \multicolumn{4}{p{0.6\linewidth}}{
        \cellcolor[rgb]{0.9,0.9,0.9}{
            \makecell[{{p{\linewidth}}}]{
                \texttt{\tiny{[P1$\rangle$GM]}}
                \texttt{SPOSTA: P, (4, 4)} \\
            }
        }
    }
    & & \\ \\

    \theutterance \stepcounter{utterance}  
    & & & \multicolumn{2}{p{0.3\linewidth}}{
        \cellcolor[rgb]{0.9,0.9,0.9}{
            \makecell[{{p{\linewidth}}}]{
                \texttt{\tiny{[GM$|$GM]}}
                \texttt{SPOSTA: P, (4, 4)} \\
            }
        }
    }
    & & \\ \\

    \theutterance \stepcounter{utterance}  
    & & & \multicolumn{2}{p{0.3\linewidth}}{
        \cellcolor[rgb]{0.9,0.9,0.9}{
            \makecell[{{p{\linewidth}}}]{
                \texttt{\tiny{[GM$|$GM]}}
                \texttt{Penalità: (4,4) non è vuoto, ma contiene '┤'.} \\
\\ 
\texttt{You have collectively accumulated 8 of 16 penalties.} \\
\texttt{Riprova!} \\
            }
        }
    }
    & & \\ \\

    \theutterance \stepcounter{utterance}  
    & & \multicolumn{4}{p{0.6\linewidth}}{
        \cellcolor[rgb]{0.9,0.9,0.9}{
            \makecell[{{p{\linewidth}}}]{
                \texttt{\tiny{[P1$\langle$GM]}}
                \texttt{Penalità: (4,4) non è vuoto, ma contiene '┤'.} \\
\\ 
\texttt{You have collectively accumulated 8 of 16 penalties.} \\
\texttt{Riprova!} \\
            }
        }
    }
    & \\ \\

    \theutterance \stepcounter{utterance}  
    & \multicolumn{4}{p{0.6\linewidth}}{
        \cellcolor[rgb]{0.9,0.9,0.9}{
            \makecell[{{p{\linewidth}}}]{
                \texttt{\tiny{[P1$\rangle$GM]}}
                \texttt{SPOSTA: P, (4, 4) no... aspetta!} \\
\\ 
\texttt{Mi dispiace, ho capito che (4,4) non è vuoto. Devo riprovare.} \\
\\ 
\texttt{COMUNICA: Mi scuso per l'errore precedente! Ho capito che (4,4) non è vuoto. Potresti suggerirmi una posizione alternativa per 'P' sulla tua griglia che sia vuota e possa aiutare ad allineare le nostre griglie?} \\
            }
        }
    }
    & & \\ \\

    \theutterance \stepcounter{utterance}  
    & & & \multicolumn{2}{p{0.3\linewidth}}{
        \cellcolor[rgb]{0.9,0.9,0.9}{
            \makecell[{{p{\linewidth}}}]{
                \texttt{\tiny{[GM$|$GM]}}
                \texttt{SPOSTA: P, (4, 4) no... aspetta!} \\
\\ 
\texttt{Mi dispiace, ho capito che (4,4) non è vuoto. Devo riprovare.} \\
\\ 
\texttt{COMUNICA: Mi scuso per l'errore precedente! Ho capito che (4,4) non è vuoto. Potresti suggerirmi una posizione alternativa per 'P' sulla tua griglia che sia vuota e possa aiutare ad allineare le nostre griglie?} \\
            }
        }
    }
    & & \\ \\

    \theutterance \stepcounter{utterance}  
    & & & \multicolumn{2}{p{0.3\linewidth}}{
        \cellcolor[rgb]{0.9,0.9,0.9}{
            \makecell[{{p{\linewidth}}}]{
                \texttt{\tiny{[GM$|$GM]}}
                \texttt{Invalid response format: SPOSTA: P, (4, 4) no... aspetta!} \\
\\ 
\texttt{Mi dispiace, ho capito che (4,4) non è vuoto. Devo riprovare.} \\
\\ 
\texttt{COMUNICA: Mi scuso per l'errore precedente! Ho capito che (4,4) non è vuoto. Potresti suggerirmi una posizione alternativa per 'P' sulla tua griglia che sia vuota e possa aiutare ad allineare le nostre griglie?} \\
            }
        }
    }
    & & \\ \\

    \theutterance \stepcounter{utterance}  
    & & \multicolumn{4}{p{0.6\linewidth}}{
        \cellcolor[rgb]{0.9,0.9,0.9}{
            \makecell[{{p{\linewidth}}}]{
                \texttt{\tiny{[P1$\langle$GM]}}
                \texttt{Penalità: Il tuo messaggio contiene più di un comando!} \\
\\ 
\texttt{Assicurati che la tua risposta contenga solo `COMUNICA: <MESSAGE>` o `SPOSTA: <OBJECT>, (<X>, <Y>)`, e nient'altro!} \\
\\ 
\texttt{You have collectively accumulated 9 of 16 penalties.} \\
\texttt{Riprova!} \\
            }
        }
    }
    & \\ \\

    \theutterance \stepcounter{utterance}  
    & \multicolumn{4}{p{0.6\linewidth}}{
        \cellcolor[rgb]{0.9,0.9,0.9}{
            \makecell[{{p{\linewidth}}}]{
                \texttt{\tiny{[P1$\rangle$GM]}}
                \texttt{COMUNICA: Mi scuso per l'errore precedente! Potresti suggerirmi una posizione alternativa per 'P' sulla tua griglia che sia vuota e possa aiutare ad allineare le nostre griglie?} \\
            }
        }
    }
    & & \\ \\

    \theutterance \stepcounter{utterance}  
    & & & \multicolumn{2}{p{0.3\linewidth}}{
        \cellcolor[rgb]{0.9,0.9,0.9}{
            \makecell[{{p{\linewidth}}}]{
                \texttt{\tiny{[GM$|$GM]}}
                \texttt{COMUNICA: Mi scuso per l'errore precedente! Potresti suggerirmi una posizione alternativa per 'P' sulla tua griglia che sia vuota e possa aiutare ad allineare le nostre griglie?} \\
            }
        }
    }
    & & \\ \\

    \theutterance \stepcounter{utterance}  
    & & \multicolumn{4}{p{0.6\linewidth}}{
        \cellcolor[rgb]{0.9,0.9,0.9}{
            \makecell[{{p{\linewidth}}}]{
                \texttt{\tiny{[GM$\rangle$P2]}}
                \texttt{Il tuo messaggio è stato inoltrato all'altro giocatore.} \\
\\ 
\\ 
\texttt{You are currently playing round 6 of maximum 28.} \\
\texttt{You have collectively accumulated 9 of 16 penalties.} \\
\texttt{L'altro giocatore ha inviato questo messaggio:} \\
\\ 
\\ 
\texttt{"Mi scuso per l'errore precedente! Potresti suggerirmi una posizione alternativa per 'P' sulla tua griglia che sia vuota e possa aiutare ad allineare le nostre griglie?"} \\
\\ 
\\ 
\texttt{Qual è il tuo prossimo comando?} \\
            }
        }
    }
    & \\ \\

    \theutterance \stepcounter{utterance}  
    & & & \multicolumn{4}{p{0.6\linewidth}}{
        \cellcolor[rgb]{0.9,0.9,0.9}{
            \makecell[{{p{\linewidth}}}]{
                \texttt{\tiny{[GM$\langle$P2]}}
                \texttt{COMUNICA: Una posizione alternativa per 'P' sulla mia griglia potrebbe essere (2,1), poiché è attualmente vuota e potrebbe essere utile per allineare le nostre griglie. Tuttavia, vorrei assicurarmi che questa posizione sia coerente con la tua strategia e le tue mosse future. Potresti confermare se questa posizione ti sembra accettabile?} \\
            }
        }
    }
     \\ \\

    \theutterance \stepcounter{utterance}  
    & & & \multicolumn{2}{p{0.3\linewidth}}{
        \cellcolor[rgb]{0.9,0.9,0.9}{
            \makecell[{{p{\linewidth}}}]{
                \texttt{\tiny{[GM$|$GM]}}
                \texttt{COMUNICA: Una posizione alternativa per 'P' sulla mia griglia potrebbe essere (2,1), poiché è attualmente vuota e potrebbe essere utile per allineare le nostre griglie. Tuttavia, vorrei assicurarmi che questa posizione sia coerente con la tua strategia e le tue mosse future. Potresti confermare se questa posizione ti sembra accettabile?} \\
            }
        }
    }
    & & \\ \\

    \theutterance \stepcounter{utterance}  
    & & \multicolumn{4}{p{0.6\linewidth}}{
        \cellcolor[rgb]{0.9,0.9,0.9}{
            \makecell[{{p{\linewidth}}}]{
                \texttt{\tiny{[P1$\langle$GM]}}
                \texttt{Il tuo messaggio è stato inoltrato all'altro giocatore.} \\
\\ 
\\ 
\texttt{You are currently playing round 6 of maximum 28.} \\
\texttt{You have collectively accumulated 9 of 16 penalties.} \\
\texttt{L'altro giocatore ha inviato questo messaggio:} \\
\\ 
\\ 
\texttt{"Una posizione alternativa per 'P' sulla mia griglia potrebbe essere (2,1), poiché è attualmente vuota e potrebbe essere utile per allineare le nostre griglie. Tuttavia, vorrei assicurarmi che questa posizione sia coerente con la tua strategia e le tue mosse future. Potresti confermare se questa posizione ti sembra accettabile?"} \\
\\ 
\\ 
\texttt{Qual è il tuo prossimo comando?} \\
            }
        }
    }
    & \\ \\

    \theutterance \stepcounter{utterance}  
    & \multicolumn{4}{p{0.6\linewidth}}{
        \cellcolor[rgb]{0.9,0.9,0.9}{
            \makecell[{{p{\linewidth}}}]{
                \texttt{\tiny{[P1$\rangle$GM]}}
                \texttt{COMUNICA: Sì, la posizione (2,1) per 'P' sulla tua griglia sembra accettabile e potrebbe essere utile per allineare le nostre griglie. Propongo di spostare 'P' in quella posizione sulla tua griglia e poi discutere le prossime mosse per allineare gli altri oggetti. Come procediamo?} \\
            }
        }
    }
    & & \\ \\

    \theutterance \stepcounter{utterance}  
    & & & \multicolumn{2}{p{0.3\linewidth}}{
        \cellcolor[rgb]{0.9,0.9,0.9}{
            \makecell[{{p{\linewidth}}}]{
                \texttt{\tiny{[GM$|$GM]}}
                \texttt{COMUNICA: Sì, la posizione (2,1) per 'P' sulla tua griglia sembra accettabile e potrebbe essere utile per allineare le nostre griglie. Propongo di spostare 'P' in quella posizione sulla tua griglia e poi discutere le prossime mosse per allineare gli altri oggetti. Come procediamo?} \\
            }
        }
    }
    & & \\ \\

    \theutterance \stepcounter{utterance}  
    & & \multicolumn{4}{p{0.6\linewidth}}{
        \cellcolor[rgb]{0.9,0.9,0.9}{
            \makecell[{{p{\linewidth}}}]{
                \texttt{\tiny{[GM$\rangle$P2]}}
                \texttt{Il tuo messaggio è stato inoltrato all'altro giocatore.} \\
\\ 
\\ 
\texttt{You are currently playing round 7 of maximum 28.} \\
\texttt{You have collectively accumulated 9 of 16 penalties.} \\
\texttt{L'altro giocatore ha inviato questo messaggio:} \\
\\ 
\\ 
\texttt{"Sì, la posizione (2,1) per 'P' sulla tua griglia sembra accettabile e potrebbe essere utile per allineare le nostre griglie. Propongo di spostare 'P' in quella posizione sulla tua griglia e poi discutere le prossime mosse per allineare gli altri oggetti. Come procediamo?"} \\
\\ 
\\ 
\texttt{Qual è il tuo prossimo comando?} \\
            }
        }
    }
    & \\ \\

    \theutterance \stepcounter{utterance}  
    & & & \multicolumn{4}{p{0.6\linewidth}}{
        \cellcolor[rgb]{0.9,0.9,0.9}{
            \makecell[{{p{\linewidth}}}]{
                \texttt{\tiny{[GM$\langle$P2]}}
                \texttt{SPOSTA: P, (2,1)} \\
            }
        }
    }
     \\ \\

    \theutterance \stepcounter{utterance}  
    & & & \multicolumn{2}{p{0.3\linewidth}}{
        \cellcolor[rgb]{0.9,0.9,0.9}{
            \makecell[{{p{\linewidth}}}]{
                \texttt{\tiny{[GM$|$GM]}}
                \texttt{SPOSTA: P, (2,1)} \\
            }
        }
    }
    & & \\ \\

    \theutterance \stepcounter{utterance}  
    & & & \multicolumn{2}{p{0.3\linewidth}}{
        \cellcolor[rgb]{0.9,0.9,0.9}{
            \makecell[{{p{\linewidth}}}]{
                \texttt{\tiny{[GM$|$GM]}}
                \texttt{Spostato 'P' in (2,1) con successo. La tua griglia aggiornata appare così:} \\
\texttt{```} \\
\texttt{ 1234567} \\
\texttt{╔═══╤═╤═╗} \\
\texttt{║◌PA│◌│◌║ 1} \\
\texttt{╟─┬─┼─┴─╢ 2} \\
\texttt{║◌│D│◌◌◌║ 3} \\
\texttt{║◌├─┤◌┌─╢ 4} \\
\texttt{╟─┘◌│◌│O║ 5} \\
\texttt{║◌◌◌└┬┴─╢ 6} \\
\texttt{║◌S◌T│◌M║ 7} \\
\texttt{╚════╧══╝} \\
\texttt{```} \\
            }
        }
    }
    & & \\ \\

    \theutterance \stepcounter{utterance}  
    & & \multicolumn{4}{p{0.6\linewidth}}{
        \cellcolor[rgb]{0.9,0.9,0.9}{
            \makecell[{{p{\linewidth}}}]{
                \texttt{\tiny{[P1$\langle$GM]}}
                \texttt{Il tuo messaggio è stato inoltrato all'altro giocatore.} \\
\\ 
\\ 
\texttt{You are currently playing round 7 of maximum 28.} \\
\texttt{You have collectively accumulated 9 of 16 penalties.} \\
\texttt{L'altro giocatore ha spostato un oggetto sulla sua griglia.} \\
\\ 
\\ 
\texttt{Qual è il tuo prossimo comando?} \\
            }
        }
    }
    & \\ \\

    \theutterance \stepcounter{utterance}  
    & \multicolumn{4}{p{0.6\linewidth}}{
        \cellcolor[rgb]{0.9,0.9,0.9}{
            \makecell[{{p{\linewidth}}}]{
                \texttt{\tiny{[P1$\rangle$GM]}}
                \texttt{SPOSTA: P, (2,1)} \\
            }
        }
    }
    & & \\ \\

    \theutterance \stepcounter{utterance}  
    & & & \multicolumn{2}{p{0.3\linewidth}}{
        \cellcolor[rgb]{0.9,0.9,0.9}{
            \makecell[{{p{\linewidth}}}]{
                \texttt{\tiny{[GM$|$GM]}}
                \texttt{SPOSTA: P, (2,1)} \\
            }
        }
    }
    & & \\ \\

    \theutterance \stepcounter{utterance}  
    & & & \multicolumn{2}{p{0.3\linewidth}}{
        \cellcolor[rgb]{0.9,0.9,0.9}{
            \makecell[{{p{\linewidth}}}]{
                \texttt{\tiny{[GM$|$GM]}}
                \texttt{Penalità: (2,1) non è vuoto, ma contiene 'S'.} \\
\\ 
\texttt{You have collectively accumulated 10 of 16 penalties.} \\
\texttt{Riprova!} \\
            }
        }
    }
    & & \\ \\

    \theutterance \stepcounter{utterance}  
    & & \multicolumn{4}{p{0.6\linewidth}}{
        \cellcolor[rgb]{0.9,0.9,0.9}{
            \makecell[{{p{\linewidth}}}]{
                \texttt{\tiny{[P1$\langle$GM]}}
                \texttt{Penalità: (2,1) non è vuoto, ma contiene 'S'.} \\
\\ 
\texttt{You have collectively accumulated 10 of 16 penalties.} \\
\texttt{Riprova!} \\
            }
        }
    }
    & \\ \\

    \theutterance \stepcounter{utterance}  
    & \multicolumn{4}{p{0.6\linewidth}}{
        \cellcolor[rgb]{0.9,0.9,0.9}{
            \makecell[{{p{\linewidth}}}]{
                \texttt{\tiny{[P1$\rangle$GM]}}
                \texttt{COMUNICA: Mi dispiace, ho capito che (2,1) non è vuoto. Potresti aggiornarmi sulla nuova disposizione degli oggetti sulla tua griglia dopo aver spostato un oggetto? Questo mi aiuterà a pianificare la mia prossima mossa.} \\
            }
        }
    }
    & & \\ \\

    \theutterance \stepcounter{utterance}  
    & & & \multicolumn{2}{p{0.3\linewidth}}{
        \cellcolor[rgb]{0.9,0.9,0.9}{
            \makecell[{{p{\linewidth}}}]{
                \texttt{\tiny{[GM$|$GM]}}
                \texttt{COMUNICA: Mi dispiace, ho capito che (2,1) non è vuoto. Potresti aggiornarmi sulla nuova disposizione degli oggetti sulla tua griglia dopo aver spostato un oggetto? Questo mi aiuterà a pianificare la mia prossima mossa.} \\
            }
        }
    }
    & & \\ \\

    \theutterance \stepcounter{utterance}  
    & & \multicolumn{4}{p{0.6\linewidth}}{
        \cellcolor[rgb]{0.9,0.9,0.9}{
            \makecell[{{p{\linewidth}}}]{
                \texttt{\tiny{[GM$\rangle$P2]}}
                \texttt{Spostato 'P' in (2,1) con successo. La tua griglia aggiornata appare così:} \\
\texttt{```} \\
\texttt{ 1234567} \\
\texttt{╔═══╤═╤═╗} \\
\texttt{║◌PA│◌│◌║ 1} \\
\texttt{╟─┬─┼─┴─╢ 2} \\
\texttt{║◌│D│◌◌◌║ 3} \\
\texttt{║◌├─┤◌┌─╢ 4} \\
\texttt{╟─┘◌│◌│O║ 5} \\
\texttt{║◌◌◌└┬┴─╢ 6} \\
\texttt{║◌S◌T│◌M║ 7} \\
\texttt{╚════╧══╝} \\
\texttt{```} \\
\texttt{You are currently playing round 8 of maximum 28.} \\
\texttt{You have collectively accumulated 10 of 16 penalties.} \\
\texttt{L'altro giocatore ha inviato questo messaggio:} \\
\\ 
\\ 
\texttt{"Mi dispiace, ho capito che (2,1) non è vuoto. Potresti aggiornarmi sulla nuova disposizione degli oggetti sulla tua griglia dopo aver spostato un oggetto? Questo mi aiuterà a pianificare la mia prossima mossa."} \\
\\ 
\\ 
\texttt{Qual è il tuo prossimo comando?} \\
            }
        }
    }
    & \\ \\

    \theutterance \stepcounter{utterance}  
    & & & \multicolumn{4}{p{0.6\linewidth}}{
        \cellcolor[rgb]{0.9,0.9,0.9}{
            \makecell[{{p{\linewidth}}}]{
                \texttt{\tiny{[GM$\langle$P2]}}
                \texttt{COMUNICA: La mia griglia aggiornata dopo lo spostamento di 'P' in (2,1) è:} \\
\texttt{```} \\
\texttt{ 1234567} \\
\texttt{╔═══╤═╤═╗} \\
\texttt{║◌PA│◌│◌║ 1} \\
\texttt{╟─┬─┼─┴─╢ 2} \\
\texttt{║◌│D│◌◌◌║ 3} \\
\texttt{║◌├─┤◌┌─╢ 4} \\
\texttt{╟─┘◌│◌│O║ 5} \\
\texttt{║◌◌◌└┬┴─╢ 6} \\
\texttt{║◌S◌T│◌M║ 7} \\
\texttt{╚════╧══╝} \\
\texttt{```} \\
            }
        }
    }
     \\ \\

    \theutterance \stepcounter{utterance}  
    & & & \multicolumn{2}{p{0.3\linewidth}}{
        \cellcolor[rgb]{0.9,0.9,0.9}{
            \makecell[{{p{\linewidth}}}]{
                \texttt{\tiny{[GM$|$GM]}}
                \texttt{COMUNICA: La mia griglia aggiornata dopo lo spostamento di 'P' in (2,1) è:} \\
\texttt{```} \\
\texttt{ 1234567} \\
\texttt{╔═══╤═╤═╗} \\
\texttt{║◌PA│◌│◌║ 1} \\
\texttt{╟─┬─┼─┴─╢ 2} \\
\texttt{║◌│D│◌◌◌║ 3} \\
\texttt{║◌├─┤◌┌─╢ 4} \\
\texttt{╟─┘◌│◌│O║ 5} \\
\texttt{║◌◌◌└┬┴─╢ 6} \\
\texttt{║◌S◌T│◌M║ 7} \\
\texttt{╚════╧══╝} \\
\texttt{```} \\
            }
        }
    }
    & & \\ \\

    \theutterance \stepcounter{utterance}  
    & & & \multicolumn{2}{p{0.3\linewidth}}{
        \cellcolor[rgb]{0.9,0.9,0.9}{
            \makecell[{{p{\linewidth}}}]{
                \texttt{\tiny{[GM$|$GM]}}
                \texttt{Invalid format: tail is not empty: '} \\
\\ 
\texttt{ 1234567} \\
\texttt{╔═══╤═╤═╗} \\
\texttt{║◌PA│◌│◌║ 1} \\
\texttt{╟─┬─┼─┴─╢ 2} \\
\texttt{║◌│D│◌◌◌║ 3} \\
\texttt{║◌├─┤◌┌─╢ 4} \\
\texttt{╟─┘◌│◌│O║ 5} \\
\texttt{║◌◌◌└┬┴─╢ 6} \\
\texttt{║◌S◌T│◌M║ 7} \\
\texttt{╚════╧══╝} \\
            }
        }
    }
    & & \\ \\

    \theutterance \stepcounter{utterance}  
    & & \multicolumn{4}{p{0.6\linewidth}}{
        \cellcolor[rgb]{0.9,0.9,0.9}{
            \makecell[{{p{\linewidth}}}]{
                \texttt{\tiny{[GM$\rangle$P2]}}
                \texttt{Penalità: Il tuo messaggio non deve contenere nulla dopo il comando!} \\
\\ 
\texttt{Assicurati che la tua risposta contenga solo `COMUNICA: <MESSAGE>` o `SPOSTA: <OBJECT>, (<X>, <Y>)`, e nient'altro!} \\
\\ 
\texttt{You have collectively accumulated 11 of 16 penalties.} \\
\texttt{Riprova!} \\
            }
        }
    }
    & \\ \\

    \theutterance \stepcounter{utterance}  
    & & & \multicolumn{4}{p{0.6\linewidth}}{
        \cellcolor[rgb]{0.9,0.9,0.9}{
            \makecell[{{p{\linewidth}}}]{
                \texttt{\tiny{[GM$\langle$P2]}}
                \texttt{COMUNICA: La mia griglia aggiornata dopo lo spostamento di 'P' in (2,1) è:} \\
\texttt{ 1234567} \\
\texttt{╔═══╤═╤═╗} \\
\texttt{║◌PA│◌│◌║ 1} \\
\texttt{╟─┬─┼─┴─╢ 2} \\
\texttt{║◌│D│◌◌◌║ 3} \\
\texttt{║◌├─┤◌┌─╢ 4} \\
\texttt{╟─┘◌│◌│O║ 5} \\
\texttt{║◌◌◌└┬┴─╢ 6} \\
\texttt{║◌S◌T│◌M║ 7} \\
\texttt{╚════╧══╝} \\
            }
        }
    }
     \\ \\

    \theutterance \stepcounter{utterance}  
    & & & \multicolumn{2}{p{0.3\linewidth}}{
        \cellcolor[rgb]{0.9,0.9,0.9}{
            \makecell[{{p{\linewidth}}}]{
                \texttt{\tiny{[GM$|$GM]}}
                \texttt{COMUNICA: La mia griglia aggiornata dopo lo spostamento di 'P' in (2,1) è:} \\
\texttt{ 1234567} \\
\texttt{╔═══╤═╤═╗} \\
\texttt{║◌PA│◌│◌║ 1} \\
\texttt{╟─┬─┼─┴─╢ 2} \\
\texttt{║◌│D│◌◌◌║ 3} \\
\texttt{║◌├─┤◌┌─╢ 4} \\
\texttt{╟─┘◌│◌│O║ 5} \\
\texttt{║◌◌◌└┬┴─╢ 6} \\
\texttt{║◌S◌T│◌M║ 7} \\
\texttt{╚════╧══╝} \\
            }
        }
    }
    & & \\ \\

    \theutterance \stepcounter{utterance}  
    & & & \multicolumn{2}{p{0.3\linewidth}}{
        \cellcolor[rgb]{0.9,0.9,0.9}{
            \makecell[{{p{\linewidth}}}]{
                \texttt{\tiny{[GM$|$GM]}}
                \texttt{Invalid format: tail is not empty: '} \\
\texttt{ 1234567} \\
\texttt{╔═══╤═╤═╗} \\
\texttt{║◌PA│◌│◌║ 1} \\
\texttt{╟─┬─┼─┴─╢ 2} \\
\texttt{║◌│D│◌◌◌║ 3} \\
\texttt{║◌├─┤◌┌─╢ 4} \\
\texttt{╟─┘◌│◌│O║ 5} \\
\texttt{║◌◌◌└┬┴─╢ 6} \\
\texttt{║◌S◌T│◌M║ 7} \\
\texttt{╚════╧══╝} \\
            }
        }
    }
    & & \\ \\

    \theutterance \stepcounter{utterance}  
    & & \multicolumn{4}{p{0.6\linewidth}}{
        \cellcolor[rgb]{0.9,0.9,0.9}{
            \makecell[{{p{\linewidth}}}]{
                \texttt{\tiny{[GM$\rangle$P2]}}
                \texttt{Penalità: Il tuo messaggio non deve contenere nulla dopo il comando!} \\
\\ 
\texttt{Assicurati che la tua risposta contenga solo `COMUNICA: <MESSAGE>` o `SPOSTA: <OBJECT>, (<X>, <Y>)`, e nient'altro!} \\
\\ 
\texttt{You have collectively accumulated 12 of 16 penalties.} \\
\texttt{Riprova!} \\
            }
        }
    }
    & \\ \\

    \theutterance \stepcounter{utterance}  
    & & & \multicolumn{4}{p{0.6\linewidth}}{
        \cellcolor[rgb]{0.9,0.9,0.9}{
            \makecell[{{p{\linewidth}}}]{
                \texttt{\tiny{[GM$\langle$P2]}}
                \texttt{COMUNICA: 1234567 ╔═══╤═╤═╗ ║◌PA│◌│◌║ 1 ╟─┬─┼─┴─╢ 2 ║◌│D│◌◌◌║ 3 ║◌├─┤◌┌─╢ 4 ╟─┘◌│◌│O║ 5 ║◌◌◌└┬┴─╢ 6 ║◌S◌T│◌M║ 7 ╚════╧══╝} \\
            }
        }
    }
     \\ \\

    \theutterance \stepcounter{utterance}  
    & & & \multicolumn{2}{p{0.3\linewidth}}{
        \cellcolor[rgb]{0.9,0.9,0.9}{
            \makecell[{{p{\linewidth}}}]{
                \texttt{\tiny{[GM$|$GM]}}
                \texttt{COMUNICA: 1234567 ╔═══╤═╤═╗ ║◌PA│◌│◌║ 1 ╟─┬─┼─┴─╢ 2 ║◌│D│◌◌◌║ 3 ║◌├─┤◌┌─╢ 4 ╟─┘◌│◌│O║ 5 ║◌◌◌└┬┴─╢ 6 ║◌S◌T│◌M║ 7 ╚════╧══╝} \\
            }
        }
    }
    & & \\ \\

    \theutterance \stepcounter{utterance}  
    & & \multicolumn{4}{p{0.6\linewidth}}{
        \cellcolor[rgb]{0.9,0.9,0.9}{
            \makecell[{{p{\linewidth}}}]{
                \texttt{\tiny{[P1$\langle$GM]}}
                \texttt{Il tuo messaggio è stato inoltrato all'altro giocatore.} \\
\\ 
\\ 
\texttt{You are currently playing round 8 of maximum 28.} \\
\texttt{You have collectively accumulated 12 of 16 penalties.} \\
\texttt{L'altro giocatore ha inviato questo messaggio:} \\
\\ 
\\ 
\texttt{"1234567 ╔═══╤═╤═╗ ║◌PA│◌│◌║ 1 ╟─┬─┼─┴─╢ 2 ║◌│D│◌◌◌║ 3 ║◌├─┤◌┌─╢ 4 ╟─┘◌│◌│O║ 5 ║◌◌◌└┬┴─╢ 6 ║◌S◌T│◌M║ 7 ╚════╧══╝"} \\
\\ 
\\ 
\texttt{Qual è il tuo prossimo comando?} \\
            }
        }
    }
    & \\ \\

    \theutterance \stepcounter{utterance}  
    & \multicolumn{4}{p{0.6\linewidth}}{
        \cellcolor[rgb]{0.9,0.9,0.9}{
            \makecell[{{p{\linewidth}}}]{
                \texttt{\tiny{[P1$\rangle$GM]}}
                \texttt{SPOSTA: D, (4,5)} \\
            }
        }
    }
    & & \\ \\

    \theutterance \stepcounter{utterance}  
    & & & \multicolumn{2}{p{0.3\linewidth}}{
        \cellcolor[rgb]{0.9,0.9,0.9}{
            \makecell[{{p{\linewidth}}}]{
                \texttt{\tiny{[GM$|$GM]}}
                \texttt{SPOSTA: D, (4,5)} \\
            }
        }
    }
    & & \\ \\

    \theutterance \stepcounter{utterance}  
    & & & \multicolumn{2}{p{0.3\linewidth}}{
        \cellcolor[rgb]{0.9,0.9,0.9}{
            \makecell[{{p{\linewidth}}}]{
                \texttt{\tiny{[GM$|$GM]}}
                \texttt{Penalità: (4,5) non è vuoto, ma contiene '│'.} \\
\\ 
\texttt{You have collectively accumulated 13 of 16 penalties.} \\
\texttt{Riprova!} \\
            }
        }
    }
    & & \\ \\

    \theutterance \stepcounter{utterance}  
    & & \multicolumn{4}{p{0.6\linewidth}}{
        \cellcolor[rgb]{0.9,0.9,0.9}{
            \makecell[{{p{\linewidth}}}]{
                \texttt{\tiny{[P1$\langle$GM]}}
                \texttt{Penalità: (4,5) non è vuoto, ma contiene '│'.} \\
\\ 
\texttt{You have collectively accumulated 13 of 16 penalties.} \\
\texttt{Riprova!} \\
            }
        }
    }
    & \\ \\

    \theutterance \stepcounter{utterance}  
    & \multicolumn{4}{p{0.6\linewidth}}{
        \cellcolor[rgb]{0.9,0.9,0.9}{
            \makecell[{{p{\linewidth}}}]{
                \texttt{\tiny{[P1$\rangle$GM]}}
                \texttt{COMUNICA: Mi scuso per l'errore! La griglia che hai condiviso mostra che 'D' è già alla posizione (3,3). Potresti suggerirmi una mossa per spostare 'D' in una posizione che sia coerente con la nostra strategia di allineare le griglie?} \\
            }
        }
    }
    & & \\ \\

    \theutterance \stepcounter{utterance}  
    & & & \multicolumn{2}{p{0.3\linewidth}}{
        \cellcolor[rgb]{0.9,0.9,0.9}{
            \makecell[{{p{\linewidth}}}]{
                \texttt{\tiny{[GM$|$GM]}}
                \texttt{COMUNICA: Mi scuso per l'errore! La griglia che hai condiviso mostra che 'D' è già alla posizione (3,3). Potresti suggerirmi una mossa per spostare 'D' in una posizione che sia coerente con la nostra strategia di allineare le griglie?} \\
            }
        }
    }
    & & \\ \\

    \theutterance \stepcounter{utterance}  
    & & \multicolumn{4}{p{0.6\linewidth}}{
        \cellcolor[rgb]{0.9,0.9,0.9}{
            \makecell[{{p{\linewidth}}}]{
                \texttt{\tiny{[GM$\rangle$P2]}}
                \texttt{Il tuo messaggio è stato inoltrato all'altro giocatore.} \\
\\ 
\\ 
\texttt{You are currently playing round 9 of maximum 28.} \\
\texttt{You have collectively accumulated 13 of 16 penalties.} \\
\texttt{L'altro giocatore ha inviato questo messaggio:} \\
\\ 
\\ 
\texttt{"Mi scuso per l'errore! La griglia che hai condiviso mostra che 'D' è già alla posizione (3,3). Potresti suggerirmi una mossa per spostare 'D' in una posizione che sia coerente con la nostra strategia di allineare le griglie?"} \\
\\ 
\\ 
\texttt{Qual è il tuo prossimo comando?} \\
            }
        }
    }
    & \\ \\

    \theutterance \stepcounter{utterance}  
    & & & \multicolumn{4}{p{0.6\linewidth}}{
        \cellcolor[rgb]{0.9,0.9,0.9}{
            \makecell[{{p{\linewidth}}}]{
                \texttt{\tiny{[GM$\langle$P2]}}
                \texttt{SPOSTA: D, (4,5)} \\
            }
        }
    }
     \\ \\

    \theutterance \stepcounter{utterance}  
    & & & \multicolumn{2}{p{0.3\linewidth}}{
        \cellcolor[rgb]{0.9,0.9,0.9}{
            \makecell[{{p{\linewidth}}}]{
                \texttt{\tiny{[GM$|$GM]}}
                \texttt{SPOSTA: D, (4,5)} \\
            }
        }
    }
    & & \\ \\

    \theutterance \stepcounter{utterance}  
    & & & \multicolumn{2}{p{0.3\linewidth}}{
        \cellcolor[rgb]{0.9,0.9,0.9}{
            \makecell[{{p{\linewidth}}}]{
                \texttt{\tiny{[GM$|$GM]}}
                \texttt{Penalità: (4,5) non è vuoto, ma contiene '│'.} \\
\\ 
\texttt{You have collectively accumulated 14 of 16 penalties.} \\
\texttt{Riprova!} \\
            }
        }
    }
    & & \\ \\

    \theutterance \stepcounter{utterance}  
    & & \multicolumn{4}{p{0.6\linewidth}}{
        \cellcolor[rgb]{0.9,0.9,0.9}{
            \makecell[{{p{\linewidth}}}]{
                \texttt{\tiny{[GM$\rangle$P2]}}
                \texttt{Penalità: (4,5) non è vuoto, ma contiene '│'.} \\
\\ 
\texttt{You have collectively accumulated 14 of 16 penalties.} \\
\texttt{Riprova!} \\
            }
        }
    }
    & \\ \\

    \theutterance \stepcounter{utterance}  
    & & & \multicolumn{4}{p{0.6\linewidth}}{
        \cellcolor[rgb]{0.9,0.9,0.9}{
            \makecell[{{p{\linewidth}}}]{
                \texttt{\tiny{[GM$\langle$P2]}}
                \texttt{SPOSTA: D, (5,3)} \\
            }
        }
    }
     \\ \\

    \theutterance \stepcounter{utterance}  
    & & & \multicolumn{2}{p{0.3\linewidth}}{
        \cellcolor[rgb]{0.9,0.9,0.9}{
            \makecell[{{p{\linewidth}}}]{
                \texttt{\tiny{[GM$|$GM]}}
                \texttt{SPOSTA: D, (5,3)} \\
            }
        }
    }
    & & \\ \\

    \theutterance \stepcounter{utterance}  
    & & & \multicolumn{2}{p{0.3\linewidth}}{
        \cellcolor[rgb]{0.9,0.9,0.9}{
            \makecell[{{p{\linewidth}}}]{
                \texttt{\tiny{[GM$|$GM]}}
                \texttt{Spostato 'D' in (5,3) con successo. La tua griglia aggiornata appare così:} \\
\texttt{```} \\
\texttt{ 1234567} \\
\texttt{╔═══╤═╤═╗} \\
\texttt{║◌PA│◌│◌║ 1} \\
\texttt{╟─┬─┼─┴─╢ 2} \\
\texttt{║◌│◌│D◌◌║ 3} \\
\texttt{║◌├─┤◌┌─╢ 4} \\
\texttt{╟─┘◌│◌│O║ 5} \\
\texttt{║◌◌◌└┬┴─╢ 6} \\
\texttt{║◌S◌T│◌M║ 7} \\
\texttt{╚════╧══╝} \\
\texttt{```} \\
            }
        }
    }
    & & \\ \\

    \theutterance \stepcounter{utterance}  
    & & \multicolumn{4}{p{0.6\linewidth}}{
        \cellcolor[rgb]{0.9,0.9,0.9}{
            \makecell[{{p{\linewidth}}}]{
                \texttt{\tiny{[P1$\langle$GM]}}
                \texttt{Il tuo messaggio è stato inoltrato all'altro giocatore.} \\
\\ 
\\ 
\texttt{You are currently playing round 9 of maximum 28.} \\
\texttt{You have collectively accumulated 14 of 16 penalties.} \\
\texttt{L'altro giocatore ha spostato un oggetto sulla sua griglia.} \\
\\ 
\\ 
\texttt{Qual è il tuo prossimo comando?} \\
            }
        }
    }
    & \\ \\

    \theutterance \stepcounter{utterance}  
    & \multicolumn{4}{p{0.6\linewidth}}{
        \cellcolor[rgb]{0.9,0.9,0.9}{
            \makecell[{{p{\linewidth}}}]{
                \texttt{\tiny{[P1$\rangle$GM]}}
                \texttt{SPOSTA: S, (1,1)} \\
            }
        }
    }
    & & \\ \\

    \theutterance \stepcounter{utterance}  
    & & & \multicolumn{2}{p{0.3\linewidth}}{
        \cellcolor[rgb]{0.9,0.9,0.9}{
            \makecell[{{p{\linewidth}}}]{
                \texttt{\tiny{[GM$|$GM]}}
                \texttt{SPOSTA: S, (1,1)} \\
            }
        }
    }
    & & \\ \\

    \theutterance \stepcounter{utterance}  
    & & & \multicolumn{2}{p{0.3\linewidth}}{
        \cellcolor[rgb]{0.9,0.9,0.9}{
            \makecell[{{p{\linewidth}}}]{
                \texttt{\tiny{[GM$|$GM]}}
                \texttt{Spostato 'S' in (1,1) con successo. La tua griglia aggiornata appare così:} \\
\texttt{```} \\
\texttt{ 1234567} \\
\texttt{╔═══╤═╤═╗} \\
\texttt{║S◌◌│◌│A║ 1} \\
\texttt{╟─┬─┼─┴─╢ 2} \\
\texttt{║◌│◌│◌◌◌║ 3} \\
\texttt{║P├─┤D┌─╢ 4} \\
\texttt{╟─┘◌│◌│M║ 5} \\
\texttt{║◌◌◌└┬┴─╢ 6} \\
\texttt{║◌OT◌│◌◌║ 7} \\
\texttt{╚════╧══╝} \\
\texttt{```} \\
            }
        }
    }
    & & \\ \\

    \theutterance \stepcounter{utterance}  
    & & \multicolumn{4}{p{0.6\linewidth}}{
        \cellcolor[rgb]{0.9,0.9,0.9}{
            \makecell[{{p{\linewidth}}}]{
                \texttt{\tiny{[GM$\rangle$P2]}}
                \texttt{Spostato 'D' in (5,3) con successo. La tua griglia aggiornata appare così:} \\
\texttt{```} \\
\texttt{ 1234567} \\
\texttt{╔═══╤═╤═╗} \\
\texttt{║◌PA│◌│◌║ 1} \\
\texttt{╟─┬─┼─┴─╢ 2} \\
\texttt{║◌│◌│D◌◌║ 3} \\
\texttt{║◌├─┤◌┌─╢ 4} \\
\texttt{╟─┘◌│◌│O║ 5} \\
\texttt{║◌◌◌└┬┴─╢ 6} \\
\texttt{║◌S◌T│◌M║ 7} \\
\texttt{╚════╧══╝} \\
\texttt{```} \\
\texttt{You are currently playing round 10 of maximum 28.} \\
\texttt{You have collectively accumulated 14 of 16 penalties.} \\
\texttt{L'altro giocatore ha spostato un oggetto sulla sua griglia.} \\
\\ 
\\ 
\texttt{Qual è il tuo prossimo comando?} \\
            }
        }
    }
    & \\ \\

    \theutterance \stepcounter{utterance}  
    & & & \multicolumn{4}{p{0.6\linewidth}}{
        \cellcolor[rgb]{0.9,0.9,0.9}{
            \makecell[{{p{\linewidth}}}]{
                \texttt{\tiny{[GM$\langle$P2]}}
                \texttt{COMUNICA: La mia griglia aggiornata è: } \\
\texttt{1234567} \\
\texttt{╔═══╤═╤═╗} \\
\texttt{║◌PA│◌│◌║ 1} \\
\texttt{╟─┬─┼─┴─╢ 2} \\
\texttt{║◌│◌│D◌◌║ 3} \\
\texttt{║◌├─┤◌┌─╢ 4} \\
\texttt{╟─┘◌│◌│O║ 5} \\
\texttt{║◌◌◌└┬┴─╢ 6} \\
\texttt{║◌S◌T│◌M║ 7} \\
\texttt{╚════╧══╝} \\
            }
        }
    }
     \\ \\

    \theutterance \stepcounter{utterance}  
    & & & \multicolumn{2}{p{0.3\linewidth}}{
        \cellcolor[rgb]{0.9,0.9,0.9}{
            \makecell[{{p{\linewidth}}}]{
                \texttt{\tiny{[GM$|$GM]}}
                \texttt{COMUNICA: La mia griglia aggiornata è: } \\
\texttt{1234567} \\
\texttt{╔═══╤═╤═╗} \\
\texttt{║◌PA│◌│◌║ 1} \\
\texttt{╟─┬─┼─┴─╢ 2} \\
\texttt{║◌│◌│D◌◌║ 3} \\
\texttt{║◌├─┤◌┌─╢ 4} \\
\texttt{╟─┘◌│◌│O║ 5} \\
\texttt{║◌◌◌└┬┴─╢ 6} \\
\texttt{║◌S◌T│◌M║ 7} \\
\texttt{╚════╧══╝} \\
            }
        }
    }
    & & \\ \\

    \theutterance \stepcounter{utterance}  
    & & & \multicolumn{2}{p{0.3\linewidth}}{
        \cellcolor[rgb]{0.9,0.9,0.9}{
            \makecell[{{p{\linewidth}}}]{
                \texttt{\tiny{[GM$|$GM]}}
                \texttt{Invalid format: tail is not empty: '} \\
\texttt{1234567} \\
\texttt{╔═══╤═╤═╗} \\
\texttt{║◌PA│◌│◌║ 1} \\
\texttt{╟─┬─┼─┴─╢ 2} \\
\texttt{║◌│◌│D◌◌║ 3} \\
\texttt{║◌├─┤◌┌─╢ 4} \\
\texttt{╟─┘◌│◌│O║ 5} \\
\texttt{║◌◌◌└┬┴─╢ 6} \\
\texttt{║◌S◌T│◌M║ 7} \\
\texttt{╚════╧══╝} \\
            }
        }
    }
    & & \\ \\

    \theutterance \stepcounter{utterance}  
    & & \multicolumn{4}{p{0.6\linewidth}}{
        \cellcolor[rgb]{0.9,0.9,0.9}{
            \makecell[{{p{\linewidth}}}]{
                \texttt{\tiny{[GM$\rangle$P2]}}
                \texttt{Penalità: Il tuo messaggio non deve contenere nulla dopo il comando!} \\
\\ 
\texttt{Assicurati che la tua risposta contenga solo `COMUNICA: <MESSAGE>` o `SPOSTA: <OBJECT>, (<X>, <Y>)`, e nient'altro!} \\
\\ 
\texttt{You have collectively accumulated 15 of 16 penalties.} \\
\texttt{Riprova!} \\
            }
        }
    }
    & \\ \\

    \theutterance \stepcounter{utterance}  
    & & & \multicolumn{4}{p{0.6\linewidth}}{
        \cellcolor[rgb]{0.9,0.9,0.9}{
            \makecell[{{p{\linewidth}}}]{
                \texttt{\tiny{[GM$\langle$P2]}}
                \texttt{COMUNICA: 1234567 ╔═══╤═╤═╗ ║◌PA│◌│◌║ 1 ╟─┬─┼─┴─╢ 2 ║◌│◌│D◌◌║ 3 ║◌├─┤◌┌─╢ 4 ╟─┘◌│◌│O║ 5 ║◌◌◌└┬┴─╢ 6 ║◌S◌T│◌M║ 7 ╚════╧══╝} \\
            }
        }
    }
     \\ \\

    \theutterance \stepcounter{utterance}  
    & & & \multicolumn{2}{p{0.3\linewidth}}{
        \cellcolor[rgb]{0.9,0.9,0.9}{
            \makecell[{{p{\linewidth}}}]{
                \texttt{\tiny{[GM$|$GM]}}
                \texttt{COMUNICA: 1234567 ╔═══╤═╤═╗ ║◌PA│◌│◌║ 1 ╟─┬─┼─┴─╢ 2 ║◌│◌│D◌◌║ 3 ║◌├─┤◌┌─╢ 4 ╟─┘◌│◌│O║ 5 ║◌◌◌└┬┴─╢ 6 ║◌S◌T│◌M║ 7 ╚════╧══╝} \\
            }
        }
    }
    & & \\ \\

    \theutterance \stepcounter{utterance}  
    & & \multicolumn{4}{p{0.6\linewidth}}{
        \cellcolor[rgb]{0.9,0.9,0.9}{
            \makecell[{{p{\linewidth}}}]{
                \texttt{\tiny{[P1$\langle$GM]}}
                \texttt{Spostato 'S' in (1,1) con successo. La tua griglia aggiornata appare così:} \\
\texttt{```} \\
\texttt{ 1234567} \\
\texttt{╔═══╤═╤═╗} \\
\texttt{║S◌◌│◌│A║ 1} \\
\texttt{╟─┬─┼─┴─╢ 2} \\
\texttt{║◌│◌│◌◌◌║ 3} \\
\texttt{║P├─┤D┌─╢ 4} \\
\texttt{╟─┘◌│◌│M║ 5} \\
\texttt{║◌◌◌└┬┴─╢ 6} \\
\texttt{║◌OT◌│◌◌║ 7} \\
\texttt{╚════╧══╝} \\
\texttt{```} \\
\texttt{You are currently playing round 10 of maximum 28.} \\
\texttt{You have collectively accumulated 15 of 16 penalties.} \\
\texttt{L'altro giocatore ha inviato questo messaggio:} \\
\\ 
\\ 
\texttt{"1234567 ╔═══╤═╤═╗ ║◌PA│◌│◌║ 1 ╟─┬─┼─┴─╢ 2 ║◌│◌│D◌◌║ 3 ║◌├─┤◌┌─╢ 4 ╟─┘◌│◌│O║ 5 ║◌◌◌└┬┴─╢ 6 ║◌S◌T│◌M║ 7 ╚════╧══╝"} \\
\\ 
\\ 
\texttt{Qual è il tuo prossimo comando?} \\
            }
        }
    }
    & \\ \\

    \theutterance \stepcounter{utterance}  
    & \multicolumn{4}{p{0.6\linewidth}}{
        \cellcolor[rgb]{0.9,0.9,0.9}{
            \makecell[{{p{\linewidth}}}]{
                \texttt{\tiny{[P1$\rangle$GM]}}
                \texttt{COMUNICA: Grazie per l'aggiornamento! Noterò che 'S' è ora alla posizione (1,1) sulla mia griglia e anche sulla tua griglia. Sembra che stiamo facendo progressi! Propongo di concentrarci sugli oggetti 'P', 'D', 'T', 'O' e 'M' per allinearli. Potresti suggerirmi come procedere con questi oggetti per allineare le nostre griglie?} \\
            }
        }
    }
    & & \\ \\

    \theutterance \stepcounter{utterance}  
    & & & \multicolumn{2}{p{0.3\linewidth}}{
        \cellcolor[rgb]{0.9,0.9,0.9}{
            \makecell[{{p{\linewidth}}}]{
                \texttt{\tiny{[GM$|$GM]}}
                \texttt{COMUNICA: Grazie per l'aggiornamento! Noterò che 'S' è ora alla posizione (1,1) sulla mia griglia e anche sulla tua griglia. Sembra che stiamo facendo progressi! Propongo di concentrarci sugli oggetti 'P', 'D', 'T', 'O' e 'M' per allinearli. Potresti suggerirmi come procedere con questi oggetti per allineare le nostre griglie?} \\
            }
        }
    }
    & & \\ \\

    \theutterance \stepcounter{utterance}  
    & & \multicolumn{4}{p{0.6\linewidth}}{
        \cellcolor[rgb]{0.9,0.9,0.9}{
            \makecell[{{p{\linewidth}}}]{
                \texttt{\tiny{[GM$\rangle$P2]}}
                \texttt{Il tuo messaggio è stato inoltrato all'altro giocatore.} \\
\\ 
\\ 
\texttt{You are currently playing round 11 of maximum 28.} \\
\texttt{You have collectively accumulated 15 of 16 penalties.} \\
\texttt{L'altro giocatore ha inviato questo messaggio:} \\
\\ 
\\ 
\texttt{"Grazie per l'aggiornamento! Noterò che 'S' è ora alla posizione (1,1) sulla mia griglia e anche sulla tua griglia. Sembra che stiamo facendo progressi! Propongo di concentrarci sugli oggetti 'P', 'D', 'T', 'O' e 'M' per allinearli. Potresti suggerirmi come procedere con questi oggetti per allineare le nostre griglie?"} \\
\\ 
\\ 
\texttt{Qual è il tuo prossimo comando?} \\
            }
        }
    }
    & \\ \\

    \theutterance \stepcounter{utterance}  
    & & & \multicolumn{4}{p{0.6\linewidth}}{
        \cellcolor[rgb]{0.9,0.9,0.9}{
            \makecell[{{p{\linewidth}}}]{
                \texttt{\tiny{[GM$\langle$P2]}}
                \texttt{COMUNICA: Sembra che stiamo facendo progressi, ma dobbiamo assicurarci di allineare correttamente tutti gli oggetti. Potremmo iniziare valutando le posizioni attuali di 'P', 'D', 'T', 'O' e 'M' su entrambe le griglie e poi pianificare le mosse necessarie per allinearli. Vorrei sapere le posizioni attuali di questi oggetti sulla tua griglia per poter confrontarle con le mie e decidere le prossime mosse.} \\
            }
        }
    }
     \\ \\

    \theutterance \stepcounter{utterance}  
    & & & \multicolumn{2}{p{0.3\linewidth}}{
        \cellcolor[rgb]{0.9,0.9,0.9}{
            \makecell[{{p{\linewidth}}}]{
                \texttt{\tiny{[GM$|$GM]}}
                \texttt{COMUNICA: Sembra che stiamo facendo progressi, ma dobbiamo assicurarci di allineare correttamente tutti gli oggetti. Potremmo iniziare valutando le posizioni attuali di 'P', 'D', 'T', 'O' e 'M' su entrambe le griglie e poi pianificare le mosse necessarie per allinearli. Vorrei sapere le posizioni attuali di questi oggetti sulla tua griglia per poter confrontarle con le mie e decidere le prossime mosse.} \\
            }
        }
    }
    & & \\ \\

    \theutterance \stepcounter{utterance}  
    & & \multicolumn{4}{p{0.6\linewidth}}{
        \cellcolor[rgb]{0.9,0.9,0.9}{
            \makecell[{{p{\linewidth}}}]{
                \texttt{\tiny{[P1$\langle$GM]}}
                \texttt{Il tuo messaggio è stato inoltrato all'altro giocatore.} \\
\\ 
\\ 
\texttt{You are currently playing round 11 of maximum 28.} \\
\texttt{You have collectively accumulated 15 of 16 penalties.} \\
\texttt{L'altro giocatore ha inviato questo messaggio:} \\
\\ 
\\ 
\texttt{"Sembra che stiamo facendo progressi, ma dobbiamo assicurarci di allineare correttamente tutti gli oggetti. Potremmo iniziare valutando le posizioni attuali di 'P', 'D', 'T', 'O' e 'M' su entrambe le griglie e poi pianificare le mosse necessarie per allinearli. Vorrei sapere le posizioni attuali di questi oggetti sulla tua griglia per poter confrontarle con le mie e decidere le prossime mosse."} \\
\\ 
\\ 
\texttt{Qual è il tuo prossimo comando?} \\
            }
        }
    }
    & \\ \\

    \theutterance \stepcounter{utterance}  
    & \multicolumn{4}{p{0.6\linewidth}}{
        \cellcolor[rgb]{0.9,0.9,0.9}{
            \makecell[{{p{\linewidth}}}]{
                \texttt{\tiny{[P1$\rangle$GM]}}
                \texttt{COMUNICA: Le posizioni attuali degli oggetti sulla mia griglia sono: 'P' in (4,4), 'D' in (4,5), 'T' in (7,7), 'O' in (7,3) e 'M' in (5,5). Spero che questo ti aiuti a confrontare le posizioni e pianificare le prossime mosse. Potresti condividere le posizioni attuali di questi oggetti sulla tua griglia in modo che possiamo decidere insieme le mosse necessarie per allinearli?} \\
            }
        }
    }
    & & \\ \\

    \theutterance \stepcounter{utterance}  
    & & & \multicolumn{2}{p{0.3\linewidth}}{
        \cellcolor[rgb]{0.9,0.9,0.9}{
            \makecell[{{p{\linewidth}}}]{
                \texttt{\tiny{[GM$|$GM]}}
                \texttt{COMUNICA: Le posizioni attuali degli oggetti sulla mia griglia sono: 'P' in (4,4), 'D' in (4,5), 'T' in (7,7), 'O' in (7,3) e 'M' in (5,5). Spero che questo ti aiuti a confrontare le posizioni e pianificare le prossime mosse. Potresti condividere le posizioni attuali di questi oggetti sulla tua griglia in modo che possiamo decidere insieme le mosse necessarie per allinearli?} \\
            }
        }
    }
    & & \\ \\

    \theutterance \stepcounter{utterance}  
    & & \multicolumn{4}{p{0.6\linewidth}}{
        \cellcolor[rgb]{0.9,0.9,0.9}{
            \makecell[{{p{\linewidth}}}]{
                \texttt{\tiny{[GM$\rangle$P2]}}
                \texttt{Il tuo messaggio è stato inoltrato all'altro giocatore.} \\
\\ 
\\ 
\texttt{You are currently playing round 12 of maximum 28.} \\
\texttt{You have collectively accumulated 15 of 16 penalties.} \\
\texttt{L'altro giocatore ha inviato questo messaggio:} \\
\\ 
\\ 
\texttt{"Le posizioni attuali degli oggetti sulla mia griglia sono: 'P' in (4,4), 'D' in (4,5), 'T' in (7,7), 'O' in (7,3) e 'M' in (5,5). Spero che questo ti aiuti a confrontare le posizioni e pianificare le prossime mosse. Potresti condividere le posizioni attuali di questi oggetti sulla tua griglia in modo che possiamo decidere insieme le mosse necessarie per allinearli?"} \\
\\ 
\\ 
\texttt{Qual è il tuo prossimo comando?} \\
            }
        }
    }
    & \\ \\

    \theutterance \stepcounter{utterance}  
    & & & \multicolumn{4}{p{0.6\linewidth}}{
        \cellcolor[rgb]{0.9,0.9,0.9}{
            \makecell[{{p{\linewidth}}}]{
                \texttt{\tiny{[GM$\langle$P2]}}
                \texttt{COMUNICA: Le posizioni attuali degli oggetti sulla mia griglia sono: 'P' in (1,1), 'D' in (5,3), 'T' in (7,7), 'O' in (5,5) e 'M' in (7,7).} \\
            }
        }
    }
     \\ \\

    \theutterance \stepcounter{utterance}  
    & & & \multicolumn{2}{p{0.3\linewidth}}{
        \cellcolor[rgb]{0.9,0.9,0.9}{
            \makecell[{{p{\linewidth}}}]{
                \texttt{\tiny{[GM$|$GM]}}
                \texttt{COMUNICA: Le posizioni attuali degli oggetti sulla mia griglia sono: 'P' in (1,1), 'D' in (5,3), 'T' in (7,7), 'O' in (5,5) e 'M' in (7,7).} \\
            }
        }
    }
    & & \\ \\

    \theutterance \stepcounter{utterance}  
    & & \multicolumn{4}{p{0.6\linewidth}}{
        \cellcolor[rgb]{0.9,0.9,0.9}{
            \makecell[{{p{\linewidth}}}]{
                \texttt{\tiny{[P1$\langle$GM]}}
                \texttt{Il tuo messaggio è stato inoltrato all'altro giocatore.} \\
\\ 
\\ 
\texttt{You are currently playing round 12 of maximum 28.} \\
\texttt{You have collectively accumulated 15 of 16 penalties.} \\
\texttt{L'altro giocatore ha inviato questo messaggio:} \\
\\ 
\\ 
\texttt{"Le posizioni attuali degli oggetti sulla mia griglia sono: 'P' in (1,1), 'D' in (5,3), 'T' in (7,7), 'O' in (5,5) e 'M' in (7,7)."} \\
\\ 
\\ 
\texttt{Qual è il tuo prossimo comando?} \\
            }
        }
    }
    & \\ \\

    \theutterance \stepcounter{utterance}  
    & \multicolumn{4}{p{0.6\linewidth}}{
        \cellcolor[rgb]{0.9,0.9,0.9}{
            \makecell[{{p{\linewidth}}}]{
                \texttt{\tiny{[P1$\rangle$GM]}}
                \texttt{COMUNICA: Grazie per aver condiviso le posizioni degli oggetti sulla tua griglia! Ora che ho le posizioni di entrambe le griglie, posso vedere che ci sono alcune differenze. Propongo di iniziare ad allineare gli oggetti 'P' e 'D' per primi, poiché sembrano essere i più facili da allineare. Potremmo anche cercare di capire come spostare 'O' e 'M' per allinearli correttamente. Vorrei discutere ulteriormente la strategia per allineare questi oggetti e assicurarmi che siamo d'accordo prima di fare qualsiasi mossa. Cosa ne pensi?} \\
            }
        }
    }
    & & \\ \\

    \theutterance \stepcounter{utterance}  
    & & & \multicolumn{2}{p{0.3\linewidth}}{
        \cellcolor[rgb]{0.9,0.9,0.9}{
            \makecell[{{p{\linewidth}}}]{
                \texttt{\tiny{[GM$|$GM]}}
                \texttt{COMUNICA: Grazie per aver condiviso le posizioni degli oggetti sulla tua griglia! Ora che ho le posizioni di entrambe le griglie, posso vedere che ci sono alcune differenze. Propongo di iniziare ad allineare gli oggetti 'P' e 'D' per primi, poiché sembrano essere i più facili da allineare. Potremmo anche cercare di capire come spostare 'O' e 'M' per allinearli correttamente. Vorrei discutere ulteriormente la strategia per allineare questi oggetti e assicurarmi che siamo d'accordo prima di fare qualsiasi mossa. Cosa ne pensi?} \\
            }
        }
    }
    & & \\ \\

    \theutterance \stepcounter{utterance}  
    & & \multicolumn{4}{p{0.6\linewidth}}{
        \cellcolor[rgb]{0.9,0.9,0.9}{
            \makecell[{{p{\linewidth}}}]{
                \texttt{\tiny{[GM$\rangle$P2]}}
                \texttt{Il tuo messaggio è stato inoltrato all'altro giocatore.} \\
\\ 
\\ 
\texttt{You are currently playing round 13 of maximum 28.} \\
\texttt{You have collectively accumulated 15 of 16 penalties.} \\
\texttt{L'altro giocatore ha inviato questo messaggio:} \\
\\ 
\\ 
\texttt{"Grazie per aver condiviso le posizioni degli oggetti sulla tua griglia! Ora che ho le posizioni di entrambe le griglie, posso vedere che ci sono alcune differenze. Propongo di iniziare ad allineare gli oggetti 'P' e 'D' per primi, poiché sembrano essere i più facili da allineare. Potremmo anche cercare di capire come spostare 'O' e 'M' per allinearli correttamente. Vorrei discutere ulteriormente la strategia per allineare questi oggetti e assicurarmi che siamo d'accordo prima di fare qualsiasi mossa. Cosa ne pensi?"} \\
\\ 
\\ 
\texttt{Qual è il tuo prossimo comando?} \\
            }
        }
    }
    & \\ \\

    \theutterance \stepcounter{utterance}  
    & & & \multicolumn{4}{p{0.6\linewidth}}{
        \cellcolor[rgb]{0.9,0.9,0.9}{
            \makecell[{{p{\linewidth}}}]{
                \texttt{\tiny{[GM$\langle$P2]}}
                \texttt{COMUNICA: Sono d'accordo di iniziare ad allineare gli oggetti 'P' e 'D' per primi, poiché sembrano essere i più facili da allineare. Per 'P', potremmo spostarlo in (4,4) sulla mia griglia per allinearla con la tua. Per 'D', potremmo spostarlo in (4,5) sulla mia griglia per allinearla con la tua. Vorrei discutere ulteriormente la strategia per allineare 'O' e 'M' e assicurarmi che siamo d'accordo prima di fare qualsiasi mossa.} \\
            }
        }
    }
     \\ \\

    \theutterance \stepcounter{utterance}  
    & & & \multicolumn{2}{p{0.3\linewidth}}{
        \cellcolor[rgb]{0.9,0.9,0.9}{
            \makecell[{{p{\linewidth}}}]{
                \texttt{\tiny{[GM$|$GM]}}
                \texttt{COMUNICA: Sono d'accordo di iniziare ad allineare gli oggetti 'P' e 'D' per primi, poiché sembrano essere i più facili da allineare. Per 'P', potremmo spostarlo in (4,4) sulla mia griglia per allinearla con la tua. Per 'D', potremmo spostarlo in (4,5) sulla mia griglia per allinearla con la tua. Vorrei discutere ulteriormente la strategia per allineare 'O' e 'M' e assicurarmi che siamo d'accordo prima di fare qualsiasi mossa.} \\
            }
        }
    }
    & & \\ \\

    \theutterance \stepcounter{utterance}  
    & & \multicolumn{4}{p{0.6\linewidth}}{
        \cellcolor[rgb]{0.9,0.9,0.9}{
            \makecell[{{p{\linewidth}}}]{
                \texttt{\tiny{[P1$\langle$GM]}}
                \texttt{Il tuo messaggio è stato inoltrato all'altro giocatore.} \\
\\ 
\\ 
\texttt{You are currently playing round 13 of maximum 28.} \\
\texttt{You have collectively accumulated 15 of 16 penalties.} \\
\texttt{L'altro giocatore ha inviato questo messaggio:} \\
\\ 
\\ 
\texttt{"Sono d'accordo di iniziare ad allineare gli oggetti 'P' e 'D' per primi, poiché sembrano essere i più facili da allineare. Per 'P', potremmo spostarlo in (4,4) sulla mia griglia per allinearla con la tua. Per 'D', potremmo spostarlo in (4,5) sulla mia griglia per allinearla con la tua. Vorrei discutere ulteriormente la strategia per allineare 'O' e 'M' e assicurarmi che siamo d'accordo prima di fare qualsiasi mossa."} \\
\\ 
\\ 
\texttt{Qual è il tuo prossimo comando?} \\
            }
        }
    }
    & \\ \\

    \theutterance \stepcounter{utterance}  
    & \multicolumn{4}{p{0.6\linewidth}}{
        \cellcolor[rgb]{0.9,0.9,0.9}{
            \makecell[{{p{\linewidth}}}]{
                \texttt{\tiny{[P1$\rangle$GM]}}
                \texttt{COMUNICA: Sono d'accordo con la tua proposta di spostare 'P' in (4,4) e 'D' in (4,5) sulla tua griglia. Per 'O' e 'M', noto che sulla mia griglia 'O' è in (7,3) e 'M' è in (5,5), mentre sulla tua griglia 'O' è in (5,5) e 'M' è in (7,7). Propongo di spostare 'O' in (7,3) sulla tua griglia e 'M' in (5,5) sulla tua griglia per allinearli con la mia. Cosa ne pensi di questa proposta? Dovremmo provare a spostarli in queste posizioni?} \\
            }
        }
    }
    & & \\ \\

    \theutterance \stepcounter{utterance}  
    & & & \multicolumn{2}{p{0.3\linewidth}}{
        \cellcolor[rgb]{0.9,0.9,0.9}{
            \makecell[{{p{\linewidth}}}]{
                \texttt{\tiny{[GM$|$GM]}}
                \texttt{COMUNICA: Sono d'accordo con la tua proposta di spostare 'P' in (4,4) e 'D' in (4,5) sulla tua griglia. Per 'O' e 'M', noto che sulla mia griglia 'O' è in (7,3) e 'M' è in (5,5), mentre sulla tua griglia 'O' è in (5,5) e 'M' è in (7,7). Propongo di spostare 'O' in (7,3) sulla tua griglia e 'M' in (5,5) sulla tua griglia per allinearli con la mia. Cosa ne pensi di questa proposta? Dovremmo provare a spostarli in queste posizioni?} \\
            }
        }
    }
    & & \\ \\

    \theutterance \stepcounter{utterance}  
    & & \multicolumn{4}{p{0.6\linewidth}}{
        \cellcolor[rgb]{0.9,0.9,0.9}{
            \makecell[{{p{\linewidth}}}]{
                \texttt{\tiny{[GM$\rangle$P2]}}
                \texttt{Il tuo messaggio è stato inoltrato all'altro giocatore.} \\
\\ 
\\ 
\texttt{You are currently playing round 14 of maximum 28.} \\
\texttt{You have collectively accumulated 15 of 16 penalties.} \\
\texttt{L'altro giocatore ha inviato questo messaggio:} \\
\\ 
\\ 
\texttt{"Sono d'accordo con la tua proposta di spostare 'P' in (4,4) e 'D' in (4,5) sulla tua griglia. Per 'O' e 'M', noto che sulla mia griglia 'O' è in (7,3) e 'M' è in (5,5), mentre sulla tua griglia 'O' è in (5,5) e 'M' è in (7,7). Propongo di spostare 'O' in (7,3) sulla tua griglia e 'M' in (5,5) sulla tua griglia per allinearli con la mia. Cosa ne pensi di questa proposta? Dovremmo provare a spostarli in queste posizioni?"} \\
\\ 
\\ 
\texttt{Qual è il tuo prossimo comando?} \\
            }
        }
    }
    & \\ \\

    \theutterance \stepcounter{utterance}  
    & & & \multicolumn{4}{p{0.6\linewidth}}{
        \cellcolor[rgb]{0.9,0.9,0.9}{
            \makecell[{{p{\linewidth}}}]{
                \texttt{\tiny{[GM$\langle$P2]}}
                \texttt{SPOSTA: P, (4,4)} \\
            }
        }
    }
     \\ \\

    \theutterance \stepcounter{utterance}  
    & & & \multicolumn{2}{p{0.3\linewidth}}{
        \cellcolor[rgb]{0.9,0.9,0.9}{
            \makecell[{{p{\linewidth}}}]{
                \texttt{\tiny{[GM$|$GM]}}
                \texttt{SPOSTA: P, (4,4)} \\
            }
        }
    }
    & & \\ \\

    \theutterance \stepcounter{utterance}  
    & & & \multicolumn{2}{p{0.3\linewidth}}{
        \cellcolor[rgb]{0.9,0.9,0.9}{
            \makecell[{{p{\linewidth}}}]{
                \texttt{\tiny{[GM$|$GM]}}
                \texttt{Penalità: (4,4) non è vuoto, ma contiene '┤'.} \\
\\ 
\texttt{You have collectively accumulated 16 of 16 penalties.} \\
\texttt{Riprova!} \\
            }
        }
    }
    & & \\ \\

    \theutterance \stepcounter{utterance}  
    & & \multicolumn{4}{p{0.6\linewidth}}{
        \cellcolor[rgb]{0.9,0.9,0.9}{
            \makecell[{{p{\linewidth}}}]{
                \texttt{\tiny{[GM$\rangle$P2]}}
                \texttt{Penalità: (4,4) non è vuoto, ma contiene '┤'.} \\
\\ 
\texttt{You have collectively accumulated 16 of 16 penalties.} \\
\texttt{Riprova!} \\
            }
        }
    }
    & \\ \\

    \theutterance \stepcounter{utterance}  
    & & & \multicolumn{4}{p{0.6\linewidth}}{
        \cellcolor[rgb]{0.9,0.9,0.9}{
            \makecell[{{p{\linewidth}}}]{
                \texttt{\tiny{[GM$\langle$P2]}}
                \texttt{COMUNICA: Sono d'accordo con la tua proposta di spostare 'O' in (7,3) e 'M' in (5,5) sulla mia griglia per allinearli con la tua. Propongo di eseguire queste mosse per vedere se riusciamo ad allineare gli oggetti.} \\
            }
        }
    }
     \\ \\

    \theutterance \stepcounter{utterance}  
    & & & \multicolumn{2}{p{0.3\linewidth}}{
        \cellcolor[rgb]{0.9,0.9,0.9}{
            \makecell[{{p{\linewidth}}}]{
                \texttt{\tiny{[GM$|$GM]}}
                \texttt{COMUNICA: Sono d'accordo con la tua proposta di spostare 'O' in (7,3) e 'M' in (5,5) sulla mia griglia per allinearli con la tua. Propongo di eseguire queste mosse per vedere se riusciamo ad allineare gli oggetti.} \\
            }
        }
    }
    & & \\ \\

    \theutterance \stepcounter{utterance}  
    & & \multicolumn{4}{p{0.6\linewidth}}{
        \cellcolor[rgb]{0.9,0.9,0.9}{
            \makecell[{{p{\linewidth}}}]{
                \texttt{\tiny{[P1$\langle$GM]}}
                \texttt{Il tuo messaggio è stato inoltrato all'altro giocatore.} \\
\\ 
\\ 
\texttt{You are currently playing round 14 of maximum 28.} \\
\texttt{You have collectively accumulated 16 of 16 penalties.} \\
\texttt{L'altro giocatore ha inviato questo messaggio:} \\
\\ 
\\ 
\texttt{"Sono d'accordo con la tua proposta di spostare 'O' in (7,3) e 'M' in (5,5) sulla mia griglia per allinearli con la tua. Propongo di eseguire queste mosse per vedere se riusciamo ad allineare gli oggetti."} \\
\\ 
\\ 
\texttt{Qual è il tuo prossimo comando?} \\
            }
        }
    }
    & \\ \\

    \theutterance \stepcounter{utterance}  
    & \multicolumn{4}{p{0.6\linewidth}}{
        \cellcolor[rgb]{0.9,0.9,0.9}{
            \makecell[{{p{\linewidth}}}]{
                \texttt{\tiny{[P1$\rangle$GM]}}
                \texttt{COMUNICA: Sono felice che siamo d'accordo sulla strategia! Tuttavia, poiché abbiamo raggiunto il massimo numero di penalità, dobbiamo essere cauti con le nostre mosse. Propongo di verificare attentamente le posizioni di destinazione per 'O' e 'M' sulla tua griglia prima di effettuare qualsiasi mossa. Inoltre, potremmo valutare se è il momento di terminare il gioco inviando un messaggio "finito?" per verificare se le nostre griglie sono allineate. Cosa ne pensi?} \\
            }
        }
    }
    & & \\ \\

    \theutterance \stepcounter{utterance}  
    & & & \multicolumn{2}{p{0.3\linewidth}}{
        \cellcolor[rgb]{0.9,0.9,0.9}{
            \makecell[{{p{\linewidth}}}]{
                \texttt{\tiny{[GM$|$GM]}}
                \texttt{COMUNICA: Sono felice che siamo d'accordo sulla strategia! Tuttavia, poiché abbiamo raggiunto il massimo numero di penalità, dobbiamo essere cauti con le nostre mosse. Propongo di verificare attentamente le posizioni di destinazione per 'O' e 'M' sulla tua griglia prima di effettuare qualsiasi mossa. Inoltre, potremmo valutare se è il momento di terminare il gioco inviando un messaggio "finito?" per verificare se le nostre griglie sono allineate. Cosa ne pensi?} \\
            }
        }
    }
    & & \\ \\

    \theutterance \stepcounter{utterance}  
    & & \multicolumn{4}{p{0.6\linewidth}}{
        \cellcolor[rgb]{0.9,0.9,0.9}{
            \makecell[{{p{\linewidth}}}]{
                \texttt{\tiny{[GM$\rangle$P2]}}
                \texttt{Il tuo messaggio è stato inoltrato all'altro giocatore.} \\
\\ 
\\ 
\texttt{You are currently playing round 15 of maximum 28.} \\
\texttt{You have collectively accumulated 16 of 16 penalties.} \\
\texttt{L'altro giocatore ha inviato questo messaggio:} \\
\\ 
\\ 
\texttt{"Sono felice che siamo d'accordo sulla strategia! Tuttavia, poiché abbiamo raggiunto il massimo numero di penalità, dobbiamo essere cauti con le nostre mosse. Propongo di verificare attentamente le posizioni di destinazione per 'O' e 'M' sulla tua griglia prima di effettuare qualsiasi mossa. Inoltre, potremmo valutare se è il momento di terminare il gioco inviando un messaggio "finito?" per verificare se le nostre griglie sono allineate. Cosa ne pensi?"} \\
\\ 
\\ 
\texttt{Qual è il tuo prossimo comando?} \\
            }
        }
    }
    & \\ \\

    \theutterance \stepcounter{utterance}  
    & & & \multicolumn{4}{p{0.6\linewidth}}{
        \cellcolor[rgb]{0.9,0.9,0.9}{
            \makecell[{{p{\linewidth}}}]{
                \texttt{\tiny{[GM$\langle$P2]}}
                \texttt{COMUNICA: Finito?} \\
            }
        }
    }
     \\ \\

    \theutterance \stepcounter{utterance}  
    & & & \multicolumn{2}{p{0.3\linewidth}}{
        \cellcolor[rgb]{0.9,0.9,0.9}{
            \makecell[{{p{\linewidth}}}]{
                \texttt{\tiny{[GM$|$GM]}}
                \texttt{COMUNICA: Finito?} \\
            }
        }
    }
    & & \\ \\

    \theutterance \stepcounter{utterance}  
    & & \multicolumn{4}{p{0.6\linewidth}}{
        \cellcolor[rgb]{0.9,0.9,0.9}{
            \makecell[{{p{\linewidth}}}]{
                \texttt{\tiny{[P1$\langle$GM]}}
                \texttt{Il tuo messaggio è stato inoltrato all'altro giocatore.} \\
\\ 
\\ 
\texttt{You are currently playing round 15 of maximum 28.} \\
\texttt{You have collectively accumulated 16 of 16 penalties.} \\
\texttt{L'altro giocatore ha inviato questo messaggio:} \\
\\ 
\\ 
\texttt{"Finito?"} \\
\\ 
\\ 
\texttt{Qual è il tuo prossimo comando?} \\
            }
        }
    }
    & \\ \\

    \theutterance \stepcounter{utterance}  
    & \multicolumn{4}{p{0.6\linewidth}}{
        \cellcolor[rgb]{0.9,0.9,0.9}{
            \makecell[{{p{\linewidth}}}]{
                \texttt{\tiny{[P1$\rangle$GM]}}
                \texttt{COMUNICA: Finito!} \\
            }
        }
    }
    & & \\ \\

    \theutterance \stepcounter{utterance}  
    & & & \multicolumn{2}{p{0.3\linewidth}}{
        \cellcolor[rgb]{0.9,0.9,0.9}{
            \makecell[{{p{\linewidth}}}]{
                \texttt{\tiny{[GM$|$GM]}}
                \texttt{COMUNICA: Finito!} \\
            }
        }
    }
    & & \\ \\

    \theutterance \stepcounter{utterance}  
    & & \multicolumn{4}{p{0.6\linewidth}}{
        \cellcolor[rgb]{0.9,0.9,0.9}{
            \makecell[{{p{\linewidth}}}]{
                \texttt{\tiny{[GM$\rangle$P2]}}
                \texttt{Il tuo messaggio è stato inoltrato all'altro giocatore.} \\
\\ 
\\ 
\texttt{You are currently playing round 16 of maximum 28.} \\
\texttt{You have collectively accumulated 16 of 16 penalties.} \\
\texttt{L'altro giocatore ha inviato questo messaggio:} \\
\\ 
\\ 
\texttt{"Finito!"} \\
\\ 
\\ 
\texttt{Qual è il tuo prossimo comando?} \\
            }
        }
    }
    & \\ \\

    \theutterance \stepcounter{utterance}  
    & & & \multicolumn{4}{p{0.6\linewidth}}{
        \cellcolor[rgb]{0.9,0.9,0.9}{
            \makecell[{{p{\linewidth}}}]{
                \texttt{\tiny{[GM$\langle$P2]}}
                \texttt{COMUNICA: Finito!} \\
            }
        }
    }
     \\ \\

    \theutterance \stepcounter{utterance}  
    & & & \multicolumn{2}{p{0.3\linewidth}}{
        \cellcolor[rgb]{0.9,0.9,0.9}{
            \makecell[{{p{\linewidth}}}]{
                \texttt{\tiny{[GM$|$GM]}}
                \texttt{COMUNICA: Finito!} \\
            }
        }
    }
    & & \\ \\

    \theutterance \stepcounter{utterance}  
    & & \multicolumn{4}{p{0.6\linewidth}}{
        \cellcolor[rgb]{0.9,0.9,0.9}{
            \makecell[{{p{\linewidth}}}]{
                \texttt{\tiny{[P1$\langle$GM]}}
                \texttt{Il tuo messaggio è stato inoltrato all'altro giocatore.} \\
\\ 
\\ 
\texttt{You are currently playing round 16 of maximum 28.} \\
\texttt{You have collectively accumulated 16 of 16 penalties.} \\
\texttt{L'altro giocatore ha inviato questo messaggio:} \\
\\ 
\\ 
\texttt{"Finito!"} \\
\\ 
\\ 
\texttt{Qual è il tuo prossimo comando?} \\
            }
        }
    }
    & \\ \\

    \theutterance \stepcounter{utterance}  
    & \multicolumn{4}{p{0.6\linewidth}}{
        \cellcolor[rgb]{0.9,0.9,0.9}{
            \makecell[{{p{\linewidth}}}]{
                \texttt{\tiny{[P1$\rangle$GM]}}
                \texttt{Non posso inviare un altro comando, poiché il gioco è terminato. Il gioco è finito! Spero di aver giocato bene e di aver raggiunto un buon punteggio nonostante le penalità accumulate. Grazie per il gioco!} \\
            }
        }
    }
    & & \\ \\

    \theutterance \stepcounter{utterance}  
    & & & \multicolumn{2}{p{0.3\linewidth}}{
        \cellcolor[rgb]{0.9,0.9,0.9}{
            \makecell[{{p{\linewidth}}}]{
                \texttt{\tiny{[GM$|$GM]}}
                \texttt{Non posso inviare un altro comando, poiché il gioco è terminato. Il gioco è finito! Spero di aver giocato bene e di aver raggiunto un buon punteggio nonostante le penalità accumulate. Grazie per il gioco!} \\
            }
        }
    }
    & & \\ \\

    \theutterance \stepcounter{utterance}  
    & & & \multicolumn{2}{p{0.3\linewidth}}{
        \cellcolor[rgb]{0.9,0.9,0.9}{
            \makecell[{{p{\linewidth}}}]{
                \texttt{\tiny{[GM$|$GM]}}
                \texttt{Invalid response format} \\
            }
        }
    }
    & & \\ \\

    \theutterance \stepcounter{utterance}  
    & & & \multicolumn{2}{p{0.3\linewidth}}{
        \cellcolor[rgb]{0.9,0.9,0.9}{
            \makecell[{{p{\linewidth}}}]{
                \texttt{\tiny{[GM$|$GM]}}
                \texttt{Maximum number of penalties exceeded} \\
            }
        }
    }
    & & \\ \\

    \theutterance \stepcounter{utterance}  
    & & & \multicolumn{2}{p{0.3\linewidth}}{
        \cellcolor[rgb]{0.9,0.9,0.9}{
            \makecell[{{p{\linewidth}}}]{
                \texttt{\tiny{[GM$|$GM]}}
                \texttt{* success: False} \\
\texttt{* lose: True} \\
\texttt{* aborted: True} \\
\texttt{{-}{-}{-}{-}{-}{-}{-}} \\
\texttt{* Shifts: 2.00} \\
\texttt{* Max Shifts: 12.00} \\
\texttt{* Min Shifts: 6.00} \\
\texttt{* End Distance Sum: 22.63} \\
\texttt{* Init Distance Sum: 23.62} \\
\texttt{* Expected Distance Sum: 29.33} \\
\texttt{* Penalties: 17.00} \\
\texttt{* Max Penalties: 16.00} \\
\texttt{* Rounds: 16.00} \\
\texttt{* Max Rounds: 28.00} \\
\texttt{* Object Count: 7.00} \\
            }
        }
    }
    & & \\ \\

    \theutterance \stepcounter{utterance}  
    & & & \multicolumn{2}{p{0.3\linewidth}}{
        \cellcolor[rgb]{0.9,0.9,0.9}{
            \makecell[{{p{\linewidth}}}]{
                \texttt{\tiny{[GM$|$GM]}}
                \texttt{Initial Grids:} \\
\texttt{```} \\
\texttt{Player 1:     Player 2:} \\
\texttt{ 1234567       1234567} \\
\texttt{╔═══╤═╤═╗     ╔═══╤═╤═╗} \\
\texttt{║◌S◌│◌│A║ 1   ║P◌A│◌│◌║ 1} \\
\texttt{╟─┬─┼─┴─╢ 2   ╟─┬─┼─┴─╢ 2} \\
\texttt{║◌│◌│◌◌◌║ 3   ║◌│D│◌◌◌║ 3} \\
\texttt{║P├─┤D┌─╢ 4   ║◌├─┤◌┌─╢ 4} \\
\texttt{╟─┘◌│◌│M║ 5   ╟─┘◌│◌│O║ 5} \\
\texttt{║◌◌◌└┬┴─╢ 6   ║◌◌◌└┬┴─╢ 6} \\
\texttt{║◌OT◌│◌◌║ 7   ║◌S◌T│◌M║ 7} \\
\texttt{╚════╧══╝     ╚════╧══╝} \\
\texttt{```} \\
\\ 
\texttt{Final Grids:} \\
\texttt{```} \\
\texttt{Player 1:     Player 2:} \\
\texttt{ 1234567       1234567} \\
\texttt{╔═══╤═╤═╗     ╔═══╤═╤═╗} \\
\texttt{║S◌◌│◌│A║ 1   ║◌PA│◌│◌║ 1} \\
\texttt{╟─┬─┼─┴─╢ 2   ╟─┬─┼─┴─╢ 2} \\
\texttt{║◌│◌│◌◌◌║ 3   ║◌│◌│D◌◌║ 3} \\
\texttt{║P├─┤D┌─╢ 4   ║◌├─┤◌┌─╢ 4} \\
\texttt{╟─┘◌│◌│M║ 5   ╟─┘◌│◌│O║ 5} \\
\texttt{║◌◌◌└┬┴─╢ 6   ║◌◌◌└┬┴─╢ 6} \\
\texttt{║◌OT◌│◌◌║ 7   ║◌S◌T│◌M║ 7} \\
\texttt{╚════╧══╝     ╚════╧══╝} \\
\texttt{```} \\
            }
        }
    }
    & & \\ \\

\end{supertabular}
}

\end{document}
